%% LyX 2.0.4 created this file.  For more info, see http://www.lyx.org/.
%% Do not edit unless you really know what you are doing.
\documentclass[oneside,english, 11pt, letter]{amsart}
\usepackage[T1]{fontenc}
\usepackage[latin9]{inputenc}
\usepackage[margin=35mm]{geometry}
%\geometry{verbose,tmargin=3cm,bmargin=3cm,lmargin=3cm,rmargin=3cm}
\usepackage{textcomp}
\usepackage{amsthm}
\usepackage{amssymb}
\usepackage{esint}
\usepackage{color}
\usepackage[dvips]{graphicx}
\usepackage{epsfig}

\newcommand{\lra}[1]{\left\langle #1 \right\rangle}
\newcommand{\biglra}[1]{\big\langle #1 \big\rangle}
\newcommand{\Biglra}[1]{\Big\langle #1 \Big\rangle}

\newcommand{\ee}{\epsilon}
\newcommand{\iy}{\infty}
\newcommand{\fr}{\frac}
\newcommand{\dd}{\partial}
\newcommand{\ddbar}{\bar{\partial}}
\newcommand{\cV}{\mathcal{V}}
\newcommand{\cE}{\mathcal{E}}
\newcommand{\cG}{\mathcal{G}}
\newcommand{\ZZ}{\mathbb{Z}}
\newcommand{\HH}{\mathbb{H}}
\newcommand{\red}{\textcolor{red}}
%\newcommand{\med}{\mathcal{M}}
%\newcommand{\medstar}{\mathcal{M}^*}
\newcommand{\potker}{\mathsf{a}}
\newcommand{\Cker}{\mathsf{K}}
\newcommand{\mass}{\mathfrak{m}}
\newcommand{\ud}{\mathrm{d}}
\newcommand{\bN}{\mathbb{N}}
\newcommand{\bZ}{\mathbb{Z}}
\newcommand{\bQ}{\mathbb{Q}}
\newcommand{\bR}{\mathbb{R}}
\newcommand{\bC}{\mathbb{C}}
\newcommand{\bH}{\mathbb{H}}
\newcommand{\EX}{\mathbb{E}}
\newcommand{\PR}{\mathbb{P}}
\newcommand{\ii}{\mathrm{i}} % other options: {\imath} {\mathfrak{i}}
\newcommand{\set}[1]{\left\{ #1 \right\}}
\newcommand{\id}{\mathrm{id}}
\newcommand{\diamed}{{\diamond/m}}
%\newcommand{\clballdiamed}[1]{\left\{ z \big| \|z\|_1 \leq #1 \right\}}
%\newcommand{\clballdiamed}[1]{\left\{ z : \|z\|_1 \leq #1 \right\}}
\newcommand{\balldiamed}[1]{B_\diamed(#1)}
\newcommand{\clballdiamed}[1]{\overline{B}_\diamed(#1)}

\renewcommand{\Im}{\operatorname{Im}}
\renewcommand{\Re}{\operatorname{Re}}
\let \le \leqslant
\let \leq \leqslant
\let \ge \geqslant
\let \geq \geqslant
\let \epsilon \varepsilon

\makeatletter
%%%%%%%%%%%%%%%%%%%%%%%%%%%%%% Textclass specific LaTeX commands.
\numberwithin{equation}{section}
\numberwithin{figure}{section}
\theoremstyle{plain}
\newtheorem{thm}{\protect\theoremname}
  \theoremstyle{remark}
  \newtheorem*{claim*}{\protect\claimname}
  \theoremstyle{plain}
  \newtheorem{prop}[thm]{\protect\propositionname}
  \theoremstyle{remark}
  \newtheorem{rem}[thm]{\protect\remarkname}
  \theoremstyle{definition}
  \newtheorem{defn}[thm]{\protect\definitionname}
  \theoremstyle{plain}
  \newtheorem{lem}[thm]{\protect\lemmaname}
\numberwithin{thm}{section}

\newtheorem*{prop*}{Proposition}
\newtheorem*{thm*}{Theorem}
\newtheorem*{lem*}{Lemma}

\makeatother

\usepackage{babel}
  \providecommand{\claimname}{Claim}
  \providecommand{\definitionname}{Definition}
  \providecommand{\lemmaname}{Lemma}
  \providecommand{\propositionname}{Proposition}
  \providecommand{\remarkname}{Remark}
	\providecommand{\theoremname}{Theorem}

\begin{document}

\title[Virasoro algebra and discrete GFF]{Lattice Representations of the Virasoro Algebra I: Discrete Gaussian Free Field}
\author{Cl\'ement Hongler, Fredrik Johansson Viklund, Kalle Kyt\"ol\"a}

\begin{abstract}
Most two-dimensional massless field theories carry representations of the Virasoro
algebra as consequences of their conformal symmetry.
Recently, conformal symmetry has been rigorously established for scaling limits of
lattice models by means of discrete complex analysis, which efficiently expresses the
integrability of these models.

In this paper we study the discrete Gaussian free field
on the square grid. We show that the lattice integrability of this model gives explicit representations
of the Virasoro algebra acting on the Gibbs measures of the model. Thus, somewhat
surprisingly, the algebraic structure of Conformal Field Theory
describing the scaling limit of the model is already present on
lattice level.
\end{abstract}
\maketitle

\tableofcontents

\section{Introduction and preliminaries}


\subsection{Statistical mechanics and Conformal Field Theory }
Physical arguments suggest that 2D lattice models at continuous phase
transitions have conformally invariant scaling limits that can be
described by Conformal Field Theories (CFTs). The 2D CFTs are exactly
solvable: they can be studied in terms of representations of the Virasoro
algebra. One can in particular obtain exact formulae for the (conjectural)
scaling limits of correlations, partition functions, and critical exponents
of a large family of 2D lattice models. See, e.g.,
\cite{belavin-polyakov-zamolodchikov-i, belavin-polyakov-zamolodchikov-ii, di-francesco-mathieu-senechal}
and the references in the latter.

While CFT has been tremendously successful and allowed for the derivation
of deep and spectacular conjectures about lattice models, it usually
lends itself only to a non-rigorous approach to statistical mechanics.
Indeed, one needs to assume that the fields of the models
have conformally invariant scaling limits and that they can be described
within the framework of certain quantum field theories. In the special
case of the Gaussian free field, the field can be studied rigorously
and representations can be found in the continuum at the level of insertions
\cite{gawedzki, kang-makarov}.

An alternative approach to conformal invariance was introduced by
Schramm, with his discovery of the SLE$_\kappa$ processes \cite{schramm-sle}. The idea here
is to consider random curves that arise in the lattice models and to
describe their scaling limits using Loewner's differential equation. This
approach has been carried out rigorously for random curves in
the following models: critical percolation \cite{smirnov-percolation}, loop-erased random
walk and uniform spanning trees \cite{lsw-lerw}, the FK-Ising model
\cite{smirnov-ising}, and the discrete Gaussian free field
\cite{schramm-sheffield-gff-i}. Besides describing scaling limits of discrete
interfaces (and so scaling limits of many natural observables), the SLE process
is amenable to calculations and SLE
techniques have been successfully applied to produce rigorous proofs
of a number of the conjectures of CFT. More systematic connections between
SLE and CFT have been studied by a number of authors, with various degrees of
rigor, see, e.g., \cite{bauer-bernard-i, werner-friedrich-i, bauer-bernard-ii,
cardy-doyon-riva, kytola-local_mgales_virasoro, kang-makarov}.

Still, we are far from having a complete correspondence and, overall, it is
fair to say that the application of CFT to statistical
mechanics remains somewhat mysterious from a mathematical perspective. One
of the difficulties is to pass to the scaling limits: one has to identify
fields on lattice level, and to assume (or prove) that they have conformally
invariant scaling limits described by a local action.
Moreover, the fields do not always admit a probabilistic interpretation.
A natural approach to improve the understanding is to try
to find precursors of various CFT objects on the lattice level, and to see
what relations they satisfy.

\subsection{Exactly solvable lattice models}
A number of two-dimensional lattice models are considered exactly
solvable. Typically, this means that relations such as the Yang-Baxter
equations or discrete holomorphicity are present on the discrete level.
This implies in particular that the partition function or the correlation
functions of these models admit exact formulae. In fact, this is how most of the rigorous results about
lattice models were derived in the 20th century. See, e.g., \cite{baxter, jimbo-miwa} and references therein.

%\red {{[}Ising, 6-vertex,
%dimers{]}}.

Recent mathematical progress has allowed for rigorous proofs that
a number of exactly solvable lattice models (in the discrete sense)
have conformally invariant scaling limits, and to rigorously derive
exact formulae for their correlation functions (predicted by exactly
solved CFTs). The central tool to establish conformal invariance
of these lattice models is their discrete holomorphicity relations,
allowing for the use of discrete complex analysis techniques,
see, e.g., \cite{kenyon-ci, smirnov-percolation, smirnov-ising, chi-ising}.

Once the (lattice) exact solvability has been used to prove conformal
invariance in the limit, one can try to use conformal invariance
to reveal the algebraic structure of the Virasoro algebra, and in
this way connect the lattice solvability with the CFT solvability.
For this, one can either connect the model with SLE or the (continuum) Gaussian
free field and use the techniques described above, or try to find
lattice precursors of CFT objects \cite{benoist-hongler, dubedat-def}.
A natural question thus arises:
\begin{itemize}
\item Is there a direct connection between the exact solvability on the discrete
level and the one on the continuum level?
\end{itemize}
We will show that, surprisingly, the answer is positive in two natural cases. On the square lattice for the discrete Gaussian free field (discrete GFF) and the critical Ising model, the
same algebraic structures that arise in the continuum as consequences
of the conformal symmetry are already present on discrete level as
consequences of discrete holomorphicity. In this paper we will give the proofs for the discrete GFF and in a companion paper those for the Ising model.

This question has been investigated in the physics literature in the case of the 8-vertex model (using the Corner transfer
matrix)  and the Ising model, see \cite{itoyama-thacker, koo-saleur} and the references in the latter. However, our results are the first in which discrete holomorphicity is connected with the Virasoro algebra and our
approach gives the exact Virasoro algebras (and not deformations thereof)
with the right physical content, that is, the correct central charge. Moreover,
the action of the operators takes place on the level of Gibbs measures, and this gives concrete insight into the role of CFT for the description of probabilistic
lattice models.

Our constructions can be modified to cover the case of the discrete GFF in a half plane with Dirichlet boundary condition. This corresponds to boundary CFT (cf., e.g., \cite{cardy-i}): besides the two commuting bulk representations, we recover in the vicinity of the boundary a representation of the Virasoro algebra with correct central charge, ``mixing'' the analytic and antianalytic sectors.


\subsection{Measure-theoretic framework}
An advantage of considering lattice models before
the scaling limit is that measure-theoretic questions become much
easier to deal with.
\begin{itemize}
\setlength{\itemsep}{0.5em}

\item We consider \emph{complex Gibbs measures} on an infinite graph $\mathcal{L}$
(typically $\mathcal{L}$ is a square grid), that is, assignments of a (finite)
complex Borel measure $\mu\big|_{\mathcal{G}}$ on $\mathbb{R}^{\mathcal{G}}$
to the finite subgraphs $\mathcal{G}\subset\mathcal{L}$ satisfying
the following compatibility condition: for $\mathcal{G}\subset\mathcal{G}'$,
integrating $\mu\big|_{\mathcal{G}'}$ over $\mathbb{R}^{\mathcal{G}'\setminus\mathcal{G}}$
gives $\mu\big|_{\mathcal{G}}$.

\item If $\mu$ and $\nu$ are complex Gibbs measures on $\mathcal{L}$,
we say that $\nu$ is \emph{a change of measure} of $\mu$ if $\nu\big|_{\mathcal{G}}$
is absolutely continuous with respect to $\mu\big|_{\mathcal{G}}$
for all finite $\mathcal{G}\subset\mathcal{L}$.
\item For a linear space of complex Gibbs measures $\mathcal{M}$, we call
a linear operator $T:\mathcal{M}\to\mathcal{M}$ such $T\mu$ is a
change of measure of $\mu$ a \emph{change of measure operator}.
\end{itemize}

Our main results are constructions of explicit families of change
of measure operators acting on the Gibbs measures of the model while
forming exact representations of the Virasoro algebra that arises in CFT.
In the statements of the results, we furthermore use the following terminology:
\begin{itemize}
% KALLE: Do we really need to define symmetric/antisymmetric here?
\item We call a complex Gibbs measure $\mu$ \emph{symmetric}
%(and antisymmetric, respectively)
if for any $\mathcal{G}$ and any Borel set $\mathcal{A}\subset\mathbb{R}^{\mathcal{G}}$,
we have $\mu\big|_{\mathcal{G}}\left(-\mathcal{A}\right)=\mu\big|_{\mathcal{G}}\left(\mathcal{A}\right)$.
We call $\mu$ \emph{antisymmetric} if
$\mu\big|_{\mathcal{G}}\left(-\mathcal{A}\right)=-\mu\big|_{\mathcal{G}}\left(\mathcal{A}\right)$.
\item A change of measure operator $T$ is said to be \emph{parity preserving} (resp.
\emph{parity reversing}) if it maps
symmetric Gibbs measures to symmetric (resp. antisymmetric) Gibbs measures and
antisymmetric Gibbs measures to antisymmetric (resp. symmetric) Gibbs measures.
\item The \emph{complex conjugate} of a complex Gibbs measure $\mu$ is the Gibbs measure
$\overline{\mu}$ determined by
$\overline{\mu}\big|_{\mathcal{G}}\left(\mathcal{A}\right)
= \overline{\mu\big|_{\mathcal{G}}\left(\mathcal{A}\right)}$
for all Borel sets $\mathcal{A}\subset\mathbb{R}^{\mathcal{G}}$.
\item The \emph{complex conjugate} of a change of measure operator $T$ is the
change of measure operator $\bar{T}$ determined by
$\bar{T} \mu = \overline{T \overline{\mu}}$. Note that $\bar{T}$ is still a
$\bC$-linear operator.
\end{itemize}


\subsection{Discrete GFF}

The central object of this paper is the discrete Gaussian free field
(GFF) which we will consider on the square grid $\ZZ^2$.
The discrete GFF on the infinite square grid $\mathbb{Z}^{2}$ is defined only up
to an additive constant, see, e.g., \cite{funaki}.
To resolve this ambiguity we choose to
consider the (whole-plane) discrete Gaussian free field \emph{pinned at $0$}, that is,
conditioned to vanish at the origin.
An alternative essentially equivalent approach would be to define a discrete gradient
Gaussian field %$\nabla_e\phi$
on the edges. %and integrate it,
%For more details, see, e.g., \cite{funaki}.

%\begin{itemize}
%\setlength{\itemsep}{0.5em}
%\item
The discrete GFF on a finite connected graph $\mathcal{G}$ with Dirichlet boundary
condition on $\partial\mathcal{G}\neq\emptyset$ is a collection of centered real Gaussian random
variables $\left(\phi(x)\right)_{x \in \mathcal{G}}$ with covariances given by the Green's kernel
of $\mathcal{G}$, that is, $\mathbb{E}\left[\phi(x) \phi(y) \right]$
is the expected number of visits to $y$ of a simple random walk $(X_t)_{t=0}^\tau$ on $\mathcal{G}$
started from $x$ and stopped at the time $\tau$ of hitting $\dd \mathcal{G}$.
See, for instance, \cite{sheffield-gff-ptrf} or \cite[Chapter 9]{janson}.
This definition also works for some infinite graphs, including a case that we will also
treat explicitly: the upper half plane square lattice $\bH = \bZ \times \bN$ with
boundary $\partial \bH = \bZ \times \set{0}$.

%\item
When $\mathcal{G}$
is a finite subgraph of $\bZ^2$, the discrete GFF $\phi$ can equivalently be defined
as a Gaussian vector of $\mathbb{R}^{\mathcal{G}}$
with density proportional to $\exp\left\{-\frac{1}{2}E(\phi)\right\}$,
where $ E(\phi) := \frac{1}{4} \sum_{x\sim y} \left( \phi(x) - \phi(y) \right)^{2} $
is the Dirichlet energy, the sum being over all pairs of adjacent vertices
$x,y\in\mathcal{G}\cup\partial\mathcal{G}$, and we set $\phi\equiv0$
on $\partial\mathcal{G}$.

%\item
Simple random walk is recurrent on the infinite square grid $\bZ^2$, and so the Green's function as defined above diverges;
an ``infra-red'' regularization is needed. We will use the massive regularization: the discrete GFF
on $\bZ^2$ with mass $\mass > 0$ is a collections of centered real Gaussian random
variables $\left(\phi_\mass(x)\right)_{x \in \bZ^2}$ with covariances given by the
following massive regularization of the Green's kernel
$ G_\mass(x,y) = \sum_{t=0}^\infty (1+\mass^2)^{-1-t} \; \PR_{X_0=x} \left[ X_t=y \right] $, % \, \big| \, X_0=x
where $(X_t)_{t\in\bN}$ is simple random walk on $\bZ^2$.

\begin{itemize}
\item The \emph{pinned discrete GFF on the infinite square grid} is a collection of centered
real Gaussian random variables $\left(\phi(x)\right)_{x \in \bZ^2}$ obtained as
the massless limit of the pinned massive discrete GFFs as
\[ \phi = \lim_{\mass \searrow 0} \Big( \phi_\mass - \phi_\mass(0) \Big) .\]
\end{itemize}

%\item
The pinned discrete GFF could also be defined as the infinite-volume limit of the
discrete GFFs conditioned to vanish at the origin, that is, the limit as $N\rightarrow \infty$ of
discrete GFF on finite subgraphs $\Lambda_N \subset \ZZ^2$ with boundary conditions
$0$ on $\dd \Lambda_N \cup \{0\}$, such that the increasing sequence of finite subgraphs
exhausts the whole lattice $\ZZ^2$.
% $\Lambda_N \uparrow \ZZ^2$.
%(a particularly convenient choice is on any finite  $R_{N}\setminus\left\{ 0\right\} $, where $R_{N}$
%is the box $\left[-N,N\right]^{2}\cap\mathbb{Z}^{2}$ with $0$ boundary
%conditions on $\left\{ 0\right\} \cup\partial R_{N}$, as $N\to\infty$.
%This gives a centered Gaussian random field $\phi=\left(\phi(z)\right)_{z\in\ZZ^2}$ on $\ZZ^2$ with covariance %kernel
%given by $\mathbb{E}\left[\phi(z) \phi(w)\right] = \potker(z) + \potker(w) - \potker(z-w)$,
%where $\potker$ is the potential kernel for random walk on $\ZZ^2$, see Section~\ref{sect:monomials}
%below and \cite{lawler-limic} for more details.
%; that is, $[\Delta \potker](z) = -\delta_{0}(z)$, where $\Delta$ is the discrete Laplacian on $\ZZ^2$
%defined below in Section~\ref{sect:monomials}.

\begin{itemize}
\item{We denote by $\mu_{\mathrm{GFF}}$ the Gibbs measure of the discrete
GFF on $\mathbb{Z}^{2}$, pinned at $0$. For a finite $\mathcal{G}=\{x_1,\ldots, x_k\} \subset \ZZ^2$ and $f:\mathbb{R}^{\mathcal{G}}  \to \mathbb{C}$, we denote by $\lra{f(\phi \big|_{\mathcal{G}})}$ the expectation $\mathbb{E} \left[ f(\phi(x_1), \ldots , \phi(x_k)) \right]$ under the measure $\mu_{\mathrm{GFF}}$.
%, assuming it exists.
(Typically we consider a \emph{correlation function} and take $f=\phi\left(x_{j_1}\right)\cdots\phi\left(x_{j_m}\right)$ for $x_{j_1}, \ldots, x_{j_m} \in \mathcal{G}$ not necessarily distinct.)
Similarly, if $\nu$ is a change of measure of
$\mu_{\mathrm{GFF}}$, we write $\lra{ f(\phi \big|_{\mathcal{G}}) }_\nu$ for the integral
of $f$ with respect to $\nu\big|_{\mathcal{G}}$ (or with respect to $\nu\big|_{\mathcal{G}'}$ for
any $\mathcal{G}' \supset \mathcal{G}$). }
\end{itemize}



\subsection{Main result}

Our main result about the discrete Gaussian free field on the full plane is the following,
the more precise version of which will be stated as Theorem \ref{Ln-GFF} in Section \ref{sugawara}.
\begin{thm*}
%\label{thm:gff-virasoro-full-plane}
Let $\mu_{\mathrm{GFF}}$ be
the Gibbs measure of the discrete Gaussian free field on $\mathbb{Z}^{2}$,
pinned at $0$.
%There exists a space $\mathcal{S}$ of symmetric changes
%of measure of $\mu_{\mathrm{GFF}}$ and explicit change of measure
%operators $L_{n}:\mathcal{S}\to\mathcal{S}$ for $n\in\mathbb{Z}$
There exists a space $\mathcal{M}$ of changes of measure of $\mu_{\mathrm{GFF}}$
and explicit parity preserving change of measure operators
$L_{n}:\mathcal{M}\to\mathcal{M}$ for $n\in\mathbb{Z}$
such that $\left(L_{n}\right)_{n\in\mathbb{Z}}$ yields a representation
of the Virasoro algebra of central charge $c=1$:
\[ \left[L_{m},L_{n}\right]
= \left(m-n\right)L_{m+n}+\frac{c}{12}\left(m^{3}-m\right) \delta_{m+n,0} \, \mathrm{Id}. \]
Each operator $L_{n}$ has a complex conjugate $\bar{L}_{n}$ and
$\left(L_{n}\right)_{n\in\mathbb{Z}}$ and $\left(\bar{L}_{n}\right)_{n\in\mathbb{Z}}$
yield two commuting representations of the Virasoro algebra with central charge $c=1$.
\end{thm*}
This result is an exact discrete analogue of the classical result
that the Gaussian free field is described by a CFT of central charge
$c=1$, %(for a mathematical statement,
see for instance \cite{gawedzki, kang-makarov}.
Through the CFT--SLE correspondence \cite{bauer-bernard-i, kang-makarov},
the parameter for the corresponding SLE$_\kappa$ is $\kappa=4$; and indeed,
the level lines of a discrete GFF with appropriate boundary conditions have been shown to converge to chordal SLE${}_4$
in the lattice size scaling limit \cite{schramm-sheffield-gff-i}.

It is interesting and quite surprising that the central charge can already be seen on lattice level.

\subsection{Further results}
Similar ideas can be carried out in other settings, too.

In particular, we shall use a Coulomb gas construction to build representations of the Virasoro algebra with other central charges, see Section~\ref{sect:coulomb-gas} for precise statements. (We remark, however, that only the representations with $c=1$ preserve the symmetry of the GFF with respect to multiplication by $-1$.)

We also construct representations for the discrete GFF in the upper half-plane with Dirichlet boundary condition, see Theorem~\ref{thm:half-plane}.




%, further discussion at the end of the paper.
%
%\paragraph*{Ising model}
%
%\paragraph*{Dimers}
%
%\paragraph*{Coulomb gas}
%
%\paragraph*{Half-plane}


\subsection{Overview of construction}

Let us give an overview of the proof of Theorem \ref{thm:gff-virasoro-full-plane}
and reduce it to a number of statements to be proven in the rest
of the paper:
\begin{itemize}
\setlength{\itemsep}{0.5em}
\item We construct discrete \emph{current modes} (Section~\ref{sect:monomials}), which
are lattice analogs of $\oint \partial \phi(z) z^{n} \ud z$, $n \in \mathbb{Z}$.
\item The discrete current modes act on field insertions by discrete contour
integrals (Lemma~\ref{prop:contour-integral-current-mode}).
\item We define a space of complex Gibbs measures $\mathcal{M}$ containing $\mu_\mathrm{GFF}$ (Definition
\ref{def:space-gibbs-measures}).
\item{The discrete current modes lift to change of measure operators
$\left(a_{n}:\mathcal{M}\to\mathcal{M}\right)_{n\in\mathbb{Z}}$
(Proposition \ref{prop:current-mode-operators}).}
\item We show that the commutation relations of the $\left(a_{n}\right)_{n\in\mathbb{Z}}$ are those of the Heisenberg algebra
(Theorem~\ref{thm:heisenberg}).
\item We define the Virasoro generators $\left(L_{n}:\mathcal{M}\to\mathcal{M}\right)_{n\in\mathbb{Z}}$
from the discrete current mode operators using the Sugawara construction
(Proposition~\ref{prop:Virasoro-mode-operators}).
\item The construction extends to the antianalytic sector, yielding
two commuting representations of the Virasoro algebra (Proposition~\ref{prop:anti-analytic-sector-an}, Theorem~\ref{prop:sugawara-construction}).
\end{itemize}
Let us denote by $\mathbb{Z}_{*}^{2}$ the dual of $\mathbb{Z}^{2}$,
by $\mathbb{Z}_{\diamond}^{2}$ the diamond graph %{[}correct name?{]}
$\mathbb{Z}^{2}\cup\mathbb{Z}_{*}^{2}$, and by $\mathbb{Z}_{m}^{2}$
the medial graph of $\mathbb{Z}^{2}$ (the set of midedges of $\mathbb{Z}^{2}$), see Figure \ref{fig: square lattice}.
The lattice representation of Theorem \ref{thm:gff-virasoro-full-plane}
relies on the construction of discrete currents and monomials that
live on these graphs.
\begin{figure}
\begin{center}
\begin{minipage}{10cm}\epsfig{file=pics-lattice_Virasoro.1, width=10cm}
\end{minipage}
\end{center}
\caption{The Figure illustrates the square lattice $\mathbb{Z}^2$ and its dual $\mathbb{Z}^2_*$.
The medial lattice $\mathbb{Z}^2_m$ consists of midpoints of edges, and
the diamond graph $\mathbb{Z}^2_\diamond$ consists of points of the square lattice and its dual.}
\label{fig: square lattice}
\end{figure}
\begin{prop*}[{Sections \ref{sect:monomials}} and \ref{sect:current-modes}]
% KALLE: What?
For $k\in\mathbb{Z}$, there exist natural discrete analogues $\left[\partial\phi\right]:\mathbb{Z}_{m}^{2}\to\mathbb{C}$,
$z^{\left[k\right]}:\mathbb{Z}_{\diamond}^{2}\cup\mathbb{Z}_{m}^{2}\to\mathbb{C}$,
and $\oint_{\left[\gamma\right]}\left[\partial\phi\right](z) z^{\left[k\right]} \ud z$
of the current $\partial\phi$, the monomial functions $z \mapsto z^{k}$,
and the contour integrals $\oint_{\gamma}\partial\phi(z) z^{k} \ud z$, respectively,
such that the following holds: If
\[ J(z):= [\dd \phi](z) \]
is the discrete current, then
\[ z^{\left[0\right]} \equiv 1 \qquad \text{ and } \qquad
 \left\langle J(z) \phi(x) \right\rangle = \frac{1}{2\pi} \left( z^{\left[-1\right]}-(z-x)^{\left[-1\right]} \right)
\]
for $z\in\mathbb{Z}_{m}^{2},x \in \mathbb{Z}^{2}$. Moreover, we have
\[
2 \ii \times \oint_{\left[\gamma\right]} \left\langle J(w) J(z) \right\rangle z^{\left[k\right]} \ud z
\; = \; k \, w^{\left[k-1\right]}, \]
for all $k\in\mathbb{Z}$, $w\in\mathbb{Z}_{m}^{2}$, whenever $[\gamma]$ is a sufficiently large closed simple
positively oriented discrete contour.
\end{prop*}
Let us now
define the space of Gibbs measures relevant to our framework.
\begin{defn}
\label{def:space-gibbs-measures}
We define $\mathcal{M}$ to be the vector space of complex Gibbs measures $\mu$ such that:
\begin{itemize}
\item{The Gibbs measure $\mu \in \mathcal{M}$ is a change of measure of $\mu_{\rm{GFF}}$;}
\item{For every finite $\mathcal{G} \subset \ZZ^2$, the Radon-Nikodym derivative of $\mu\big|_{\mathcal{G}}$ with respect
to $\mu_{\rm{GFF}}\big|_{\mathcal{G}}$, denoted by $g_{\mathcal{G}}^{\mu}$,
lies in $L^p(\mu_{\rm{GFF}}\big|_{\mathcal{G}})$ for every $p < \infty$.}
\end{itemize}
We also define the subspaces $\mathcal{S} \subset \mathcal{M}$ and
$\mathcal{A} \subset \mathcal{M}$ of symmetric and antisymmetric measures
in $\mathcal{M}$, respectively. In other words, $\mathcal{S}$ and $\mathcal{A}$
consist of the measures in $\mathcal{M}$ with even and odd Radon-Nikodym derivatives
$g_{\mathcal{G}}^{\mu}$, respectively.
\end{defn}
%\begin{rem}
%The space $\mathcal{S}$ is the natural stable ``physical'' subspace of $\mathcal{M}$ that contains $\mu_{\rm{GFF}}$.
%\end{rem}

\begin{rem}\label{rem:denseness-of-polynomials}
For each $\mathcal{G}$, the linear span of polynomials in the Gaussian
variables $\{\phi(z): \, z \in \mathcal{G}\},$ is dense in
$L^p\left(\mu_{\rm{GFF}} \big|_{\mathcal{G}} \right)$ for every $0 \le p < \infty$,
see, e.g., \cite[Theorem~2.11]{janson}.
%A consequence is that it will be enough to check commutation relations for insertions
%of the current modes into polynomial correlations in Section~\ref{sect:current-modes}.
\end{rem}

\begin{lem*}[Lemma~\ref{prop:contour-integral-current-mode}]For any (not necessarily distinct)
$x_{1},\ldots,x_{m}\in\mathbb{Z}^{2}$ and every $n\in\mathbb{Z}$ the discrete
current mode $\mathfrak{a}_n$ acting on field insertions by the contour integral
\[ \lra{ \mathfrak{a}_n \phi\left(x_{1}\right)\cdots\phi\left(x_{m}\right)}
 = \frac{1}{\sqrt{\pi}}\oint_{\left[\gamma\right]}\left\langle J\left(z\right)\phi\left(x_{1}\right)\cdots\phi\left(x_{m}\right)\right\rangle z^{\left[n\right]}\mathrm{d}z \]
is independent of the choice of the contour $[\gamma]$ provided
that $[\gamma]$ encircles a large enough neighborhood of origin in the positive direction.
%it separates $\{0, x_1, \ldots, x_k\}$ from $\infty$.

More generally, the product of discrete current modes acting on field insertions by contour integrals
\begin{multline*}
\lra{\mathfrak{a}_{n_j}\cdots\mathfrak{a}_{n_1} \phi\left(x_{1}\right)\cdots\phi\left(x_{m}\right) }\\
:= \frac{1}{\pi^{j/2}}\oint_{\left[\gamma_{j}\right]} \cdots \oint_{\left[\gamma_{1}\right]} \lra{ J(w_j) \cdots J(w_1)\phi\left(x_{1}\right)\cdots\phi\left(x_{m}\right)} w_1^{\left[n_1\right]} \cdots w_j^{\left[n_j\right]} dw_1 \cdots dw_j
\end{multline*}
is independent on the choice of sequence of discrete contours $[\gamma_1] , [\gamma_2], \ldots, [\gamma_j]$
provided the contours encircle sufficiently large neighborhoods of the origin and are radially ordered:
for $l<k$ the contour $[\gamma_l]$ (corresponding to $\mathfrak{a}_{n_l}$)
is contained in the set encircled by the contour $[\gamma_k]$ (corresponding to $\mathfrak{a}_{n_k}$).
\end{lem*}
%We will always assume that integration contours corresponding to discrete current modes are radially ordered.
\begin{rem}
For any given realization, the contour integral
$\oint_{\left[\gamma\right]}J\left(z\right)z^{[n]}\mathrm{d} z$
does depend on the contour $[\gamma]$, but the expected value against any insertion
$\phi\left(x_{1}\right)\cdots\phi\left(x_{n}\right)$ is the same
for all positively oriented discrete contours $[\gamma]$ enclosing
a large enough neighborhood of the origin.
\end{rem}

By Lemma~\ref{prop:contour-integral-current-mode}, the insertions
of $\oint_{\left[\gamma\right]}J\left(z\right)z^{[n]}d z$ (and the corresponding products)
into correlation functions of the discrete GFF (and its changes of measure) are well-defined provided
we take a large enough contour $[\gamma]$ (or radially ordered sequence of contours). We show that the
current modes $(\mathfrak{a}_n)_{n\in \ZZ}$ (and their products) lift to unique change of measure
operators $\left(a_{n}\right)_{n\in\mathbb{Z}}$ which act on $\mathcal{M}$ and which can be though of as formal adjoint operators to the $\mathfrak{a}_n$. (This point of view is behind the use of $\mathfrak{a}_{-n}$ in the definition of $a_n$ below.) These change of
measure operators are the building blocks used in the construction
of the operators $\left(L_{n}\right)_{n\in\mathbb{Z}}$ of Theorem
\ref{thm:gff-virasoro-full-plane}.
\begin{prop*}[Proposition~\ref{prop:current-mode-operators}]
For each $n\in\mathbb{Z}$, there
exists a unique operator $a_{-n}:\mathcal{M}\to\mathcal{M}$
%, with $a_{-n}\mathcal{S}\subset\mathcal{A}$ and $a_{-n}\mathcal{A}\subset\mathcal{S}$,
such that for any $\mu\in\mathcal{M}$ and any
(not necessarily distinct) $x_{1},\ldots,x_{k}\in\mathbb{Z}^{2}$ and any finite $\mathcal{G}\subset\bZ^2$ containing the points,
we have
\begin{align*}
\left\langle \phi\left(x_{1}\right)\cdots\phi\left(x_{k}\right)\right\rangle _{a_{-n}\mu}
& = \lra{\mathfrak{a}_{n} \phi\left(x_{1}\right)\cdots\phi\left(x_{k}\right)}_\mu \\
& :=\frac{1}{\sqrt{\pi}}\oint_{\left[\gamma\right]}\left\langle J\left(z\right)\phi\left(x_{1}\right)\cdots\phi\left(x_{m}\right)
    g (\phi) \right\rangle z^{\left[n\right]} \mathrm{d}z,
\end{align*}
where $g$ is the Radon-Nikodym derivative of $\mu\big|_{\mathcal{G}}$ with respect
to $\mu_{\rm{GFF}}\big|_{\mathcal{G}}$, and $\left[\gamma\right]$ is any discrete
contour that encircles a large enough neighborhood of the origin.
The operator $a_{-n}$ is parity reversing.
\end{prop*}
We then show that the operators $\left(a_{n}\right)_{n\in\mathbb{Z}}$
satisfy the commutation relations of the Heisenberg algebra:
\begin{thm*}[{Theorem~\ref{thm:heisenberg}}]
\label{prop:commutation-relations}The operators $\left(a_{n}\right)_{n\in\mathbb{Z}}$
satisfy the commutation relations
\begin{eqnarray*}
\left[a_{m},a_{n}\right] & = & m\delta_{n+m,0} \, \id_\mathcal{M}
\end{eqnarray*}

\end{thm*}
Finally, we can rely on the classical Sugawara construction \cite{sugawara, sommerfield} to define
Virasoro operators $L_{n}$. See \cite{mickelsson} for a nice textbook about the algebra of the current modes. In the following statement and in the rest of the paper, a function $f : \mathbb{R}^{\ZZ^2} \to \mathbb{C}$ is said to be a \emph{cylinder function}
if it depends only on finitely many coordinates, called its \emph{base}.
\begin{thm*}[Theorem~\ref{prop:sugawara-construction}]
For each $n\in\mathbb{Z}$, in the formally infinite sum
\[ \mathfrak{L}_{n} := \frac{1}{2}\sum_{k\geq0}\mathfrak{a}_{n-k}\mathfrak{a}_{k}
    + \frac{1}{2} \sum_{k<0}\mathfrak{a}_{k}\mathfrak{a}_{n-k} ,\]
only finitely many terms are non-zero when acting on an insertion of a cylinder function.
The action of $\mathfrak{L}_{-n}$ on insertions lifts to a
well-defined parity preserving change of measure
operator $L_n \colon \mathcal{M} \to \mathcal{M}$.

The $\left(L_{n}\right)_{n\in\mathbb{Z}}$ yield a representation of
the Virasoro algebra with central charge $c=1$:
\[
\left[L_{m},L_{n}\right]=\left(m-n\right)L_{m+n}+\frac{1}{12}(m^3-m)\delta_{m+n,0} \, \id_\mathcal{M}.
\]
\end{thm*}
\begin{rem}
The above expression for Virasoro generators as a formally infinite sum of quadratic polynomials in
the currents is standard in Sugawara constructions. In usual algebraic Sugawara constructions,
the expression has well defined action in representations of an appropriate type, owing
to its truncation to only finitely many terms when applied to any particular vector ---
it is sufficient that the representation has a grading by $L_0$-eigenvalues that are bounded
from below. Our action on insertions by
$\mathfrak{L}_{n}$ truncates in a similar manner, now because of the finite base of the
inserted cylinder function. We remark, however, that the action on interesting Gibbs measures
by $L_n$ can not be recovered as a finite linear combination of quadratic expressions
in $(a_m)_{m\in\bZ}$.
\end{rem}
The operators above build a discrete version of the analytic sector
of the CFT of the free field. (One could perhaps call it
a \emph{discrete analytic sector}.) As in the continuum case, we can also
build the antianalytic sector in the same manner, and we thus obtain two
commuting representations of the Virasoro algebra.

\begin{prop*}[Proposition~\ref{prop:anti-analytic-sector-an}, Theorem~\ref{prop:sugawara-construction}]
%For each $n\in\mathbb{Z}$, there
%exists a unique operator $\bar{a}_{n}:\mathcal{M}\to\mathcal{M}$
%such that
%\[ \left\langle \phi\left(x_{1}\right)\cdots\phi\left(x_{m}\right) \right\rangle_{\bar{a}_{n}\mu}
%=\overline{\left\langle \phi\left(x_{1}\right)\cdots\phi\left(x_{m}\right)\right\rangle }_{a_{n}\overline{\mu}}
%\qquad \forall\mu\in\mathcal{M}, \, \forall x_{1},\ldots,x_{m}\in\mathbb{Z}^{2} . \]
For $n\in\bZ$, let $\bar{a}_n \colon \mathcal{M} \rightarrow \mathcal{M}$ be the
complex conjugate of the change of measure operator $a_n$.
For all $n,m\in\mathbb{Z},$ we have $\left[\bar{a}_{m},a_{n}\right]=0$
and $\left[\bar{a}_{m},\bar{a}_{n}\right]=m\delta_{m+n,0} \id_\mathcal{M}$.
Defining $L_{n}$ as in Proposition~\ref{prop:sugawara-construction}
and $\bar{L}_{n}$ as its complex conjugate,
we have that $\left(L_{n}\right)_{n\in\mathbb{Z}}$ %$\left(L_{n}:\mathcal{S}\to\mathcal{S}\right)_{n\in\mathbb{Z}}$
and $\left(\bar{L}_{n}\right)_{n\in\mathbb{Z}}$ %$\left(\bar{L}_{n}:\mathcal{S}\to\mathcal{S}\right)_{n\in\mathbb{Z}}$
yield two commuting representations of the Virasoro algebra with central
charge $c=1$:
\begin{eqnarray*}
\left[L_{n},L_{m}\right] & = & \left(n-m\right)L_{n+m}+\frac{1}{12}(m^3-m)\delta_{n+m} \, \id_\mathcal{M} \\
\left[\bar{L}_{n},\bar{L}_{m}\right] & = & \left(n-m\right)\bar{L}_{n+m}+\frac{1}{12}(m^3-m)\delta_{n+m} \, \id_\mathcal{M} \\
\left[L_{n},\bar{L}_{m}\right] & = & 0.
\end{eqnarray*}
Both representations consist of parity preserving change of measure operators.
\end{prop*}



Finally, we can use the background charge idea in the Coulomb gas construction to
yield representations of the Virasoro algebra at arbitrary central charges.
When $c \neq 1$, however, the operators $L_n^b, \bar{L}_n^b$ are
not parity preserving, and thus one loses a symmetry of the original problem.
\begin{prop*}[Proposition~\ref{prop:coulomb_gas}]
Let $c \leq 1$ and let $b = \pm \sqrt{\frac{1-c}{12}}$.
If we define the change
of measure operators $L_{n}^{b},\bar{L}_{n}^{b}:\mathcal{M}\to\mathcal{M}$
by $L_n^b:=L_n+b(n+1) a_n$, with $L_n$ and $a_n$ as above, we have that
$\left(L_{n}^{b}\right)_{n\in\mathbb{Z}}$ and
$\left(\bar{L}_{n}^{b}\right)_{n\in\mathbb{Z}}$ yield two commuting
representations of the Virasoro algebra with central charge $c$.
\end{prop*}
\subsection{General remarks}
%\red{{[}In CFT, the space on which one acts are insertions, not necessarily
%probabilistic fields{]}}

In conformal field theory, the basis $(L_n)_{n \in \ZZ}$ that spans the Virasoro
algebra is given by modes of the (analytic component of the) stress-energy
tensor $T(z)$. By comparing with the continuum GFF one might therefore attempt
to define discrete $L_n$ modes as discrete modes of (a constant times)
$[T(z)]=[J(z)]^2$. However, there seems to be no natural way to do this so
that one has independence on the choice of integration contour. Instead,
the Sugawara construction of the $L_n$ modes from the discrete current
modes can be viewed as taking place directly in ``Fourier space'' on the
level of operator valued Laurent modes. Our constructions give concrete
meaning to these objects.

%{[The role of complex measures in statistical mechanics -- Ising, etc...]}

%\red{[}Asymptotics{]}

\subsection*{Acknowledgments}
CH is supported by the Minerva Foundation and by the National Science Foundation grant number DMS-1106588.
FJV is supported by the Simons foundation.
KK is supported by the Academy of Finland.
We would like to thank St\'ephane Benoist, Dmitry Chelkak, Julien Dub\'edat,
Krzysztof Gaw{\c e}dzki, Igor Krichever, Antti Kupiainen, Jouko Mickelsson, Andrei Okounkov,
Eveliina Peltola, Stanislav Smirnov and Valerio Toledano Laredo for
interesting discussions.


\section{Discrete currents, polynomials, and contour integrals}\label{sect:monomials}
We now describe in detail the basic building blocks in our construction.

\subsection{Finite difference operators}Recall the definitions of the square grid $\ZZ^2$ with vertices $\{m+in: m,n \in \ZZ \}$, its dual $\ZZ^2_*$ with vertices $\{m+\frac{1}{2} + i(n + \frac{1}{2}): m,n \in \ZZ\}$, the diamond graph with vertices $\ZZ_{\diamond}^2 = \ZZ^2 \cup \ZZ^2_*$, and the medial graph $\ZZ^2_m$ with a vertex at the midpoint of every edge of $\ZZ^2$. We will sometimes identify an edge of a graph with its midpoint.

For $f \colon \mathbb{Z}_{\diamond}^{2}\to\mathbb{C}$ (and $f \colon \mathbb{Z}_{m}^{2}\to\mathbb{C}$, respectively), we define $\left[\partial f\right] \colon \mathbb{Z}_{m}^{2} \to \mathbb{C}$
(and $[\partial f] \colon \mathbb{Z}_{\diamond}^{2}\to\mathbb{C}$, respectively) by
\[ [\partial f](z) = \sum_{a\in\left\{ \pm\frac{1}{2},\pm\frac{\ii}{2}\right\} }\overline{a} \, f(z+a).\]
Similarly, we define
$[\bar{\partial}f ]$ by
\[ [\bar{\partial}f](z) = \sum_{a\in\left\{ \pm\frac{1}{2},\pm\frac{\ii}{2}\right\} } a \, f(z+a).\]

We say that $f \colon \mathbb{Z}_{\diamond}^{2}\to\mathbb{C}$ (or $f \colon \mathbb{Z}_{m}^{2}\to\mathbb{C}$) is \emph{discrete holomorphic} at $z$ if $[\bar{\partial} f](z) = 0$.

For a function $f$ defined on the square grid $\ZZ^2$ (more generally, on a shifted square grid,
e.g., the dual $\ZZ^2_*$), the discrete Laplacian
$[\Delta f]\colon \ZZ^2 \rightarrow \bC$ is defined by
\[ [\Delta f](z) = \frac{1}{4} \sum_{\xi \in \{\pm 1 , \pm \ii\}} f(z+\xi) - f(z), \]
where the latter sum is over the four nearest neighbors $w$ of $z$.
Note that for $f$ defined on $\ZZ^2_\diamond$ (or on $\ZZ^2_m$), the composition
$\dd \ddbar = \ddbar \dd = \Delta$ is the above Laplacian, which in particular
splits to two blocks corresponding to $\ZZ^2$ and $\ZZ^2_*$ (on $\ZZ^2_m$ the
two blocks correspond to midpoints of horizontal and vertical edges).
\begin{figure}
\begin{center}
\begin{minipage}{14cm}
\epsfig{file=pics-lattice_Virasoro.5, width=3.0cm} $\qquad$
\epsfig{file=pics-lattice_Virasoro.6, width=3.0cm} $\qquad$
\epsfig{file=pics-lattice_Virasoro.7, width=3.75cm}
\end{minipage}
\end{center}
\caption{Illustrations of the weights in the finite difference operators $\partial$ and $\bar{\partial}$ and $\Delta = \partial \bar{\partial}$ on particular sublattices. The weights of other sublattices are translations of the ones shown.}
\label{fig: discrete contours}
\end{figure}

%We also occasionally need the directional holomorphic and antiholomorphic derivatives
%in horizontal and vertical directions,
%\[ \dd^\mathrm{hor} = \frac{1}{2}(\dd + \ddbar) = \ddbar^\mathrm{hor} , \qquad
%\dd^\mathrm{ver} = \frac{1}{2}(\dd - \ddbar) = - \ddbar^\mathrm{ver} . \]
%%(and similarly for $\ddbar^\mathrm{hor}$, $\ddbar^\mathrm{ver}$).

We record a property that is used in some calculations below.
\begin{lem}\label{lem: transpose of dee}
The transposes of the finite difference operators $\dd$
are $- \dd$ (the operator $\bC^{\ZZ^2_m} \rightarrow \bC^{\ZZ^2_\diamond}$ has transpose
$\bC^{\ZZ^2_\diamond} \rightarrow \bC^{\ZZ^2_m}$ and vice versa).
In particular, if
$f \colon \ZZ^2_m \rightarrow \bC$ and $g \colon \ZZ^2_\diamond \rightarrow \bC$
are functions, at least one of which has finite support, then
\[ \sum_{z \in \ZZ^2_m} f(z) \, [\dd g](z) = - \sum_{z \in \ZZ^2_\diamond} [\dd f](z) \, g(z) . \]
Similarly, the transposes of $\ddbar$ are $-\ddbar$.
%, the transposes of $\dd^\mathrm{hor}$ are $- \dd^\mathrm{hor}$,
%and the transposes of $\dd^\mathrm{ver}$ are $- \dd^\mathrm{ver}$.
\end{lem}
%\begin{proof}
%:D
%\end{proof}


\subsection{Potential kernel and discrete Cauchy kernel}
In the full plane, the expected number of visits to a point by simple
random walk $(X_t)_{t\in\bN}$ is infinite, and for this reason a naive definition of the (massless) discrete
Gaussian free field will not work and one needs to consider a regularized version instead.

We shall use the massive regularization based on the massive Green's function $G_\mass \colon \ZZ^2 \to \bR$, $\mass > 0$,
%\[ G_\mass(z) = \EX \left[ \sum_{t=0}^\infty (1+\mass^2)^{-t-1} \mathbf{1}_{\{X_t=z\}} \right] ,
%%\qquad \text{($(X_t)_{t\in\bN}$ simple random walk on $\ZZ^2$ started at origin)}
%\]
which is the function decaying at infinity that satisfies
$\mass^2 G_\mass - [\Delta G_\mass] = \delta_0$.
%It is easy to write the Fourier transform of $G_\mass$,
%\[ G_\mass (x + \ii y) = \frac{1}{(2\pi)^2} \iint_{[-\pi,\pi]^2}
%\frac{e^{\ii (x q + y p)}}{\mass^2 + 1 - \frac{1}{2} \cos(q) - \frac{1}{2} \cos(p)} \ud q \ud p .\]
If we subtract a suitable diverging constant, for example the value at the origin,
the massless limit $\mass \searrow 0$ can be taken.
We define the  \emph{potential kernel} $\potker \colon \ZZ^2 \to \bR$ by
\begin{align*}
\potker(z)
= \; & \lim_{\mass \searrow 0} \left( G_\mass(0) - G_\mass(z) \right) . %\\ = \; &
%= \frac{1}{(2\pi)^2} \iint_{[-\pi,\pi]^2}
%\frac{1 - e^{\ii (x q + y p)}}{1 - \frac{1}{2} \cos(q) - \frac{1}{2} \cos(p)} \ud q \ud p .
\end{align*}
%The potential kernel $\potker \colon \ZZ^2 \rightarrow \bR$
It is a function satisfying $[\Delta \potker] (z) = \delta_0(z)$ and having
logarithmic growth at infinity, and it is determined up to an additive constant by these
conditions. In fact we have
\[ \potker(z) = \frac{2}{\pi} \, \log |z| + \frac{2\gamma + \log 8}{\pi} +  O(|z|^{-2})\]
as $z \rightarrow \infty$, see, e.g., \cite[Sections 4 and 6]{lawler-limic}.
For convenience, we extend the potential kernel to the
diamond lattice $\ZZ^2_\diamond = \ZZ^2 \cup \ZZ^2_*$ %,
%\[ \potker \colon \ZZ^2_\diamond \to \bR ,\]
by setting the values on $\ZZ^2_*$ to zero.
With this extension, we can consider the function
\[ [\dd \potker] \colon \ZZ^2_m \rightarrow \bC , \]
which is purely real on the horizontal edges and purely imaginary on the vertical edges of $\ZZ^2$.

See \cite{hongler-i, kenyon-ci} for discussion and references to the following result.
\begin{lem}
The function $\Cker = [\dd \potker]$ is the unique function $\Cker:\mathbb{Z}_{m}^{2}\to\mathbb{C}$
satisfying $\left[\ddbar \Cker \right](z)=\delta_{0}(z)$
and $\Cker\left(z\right) \to 0$ as $|z| \to \infty$. We
call $\Cker$ the \emph{discrete Cauchy kernel} with pole at $0$.
\end{lem}
\begin{proof}
Clearly $\Cker=[\dd \potker]$ satisfies $[\ddbar \Cker] = [\ddbar \dd \potker] = [\Delta \potker] = \delta_0$,
and $\Cker$ decays at infinity.

To show uniqueness, consider the difference of two functions satisfying the properties.
The difference is a function on $\ZZ^2_\diamond = \ZZ^2 \cup \ZZ^2_*$,
which is annihilated by $\ddbar$, and in particular its restrictions
to $\ZZ^2$ and $\ZZ^2_*$ are harmonic %, since $\Delta = \dd \ddbar$. As harmonic function
and decaying at infinity, thus zero. %, the difference is zero.
\end{proof}

\subsection{Discrete Laurent monomials}
We next define discrete Laurent monomial functions on both
$\mathbb{Z}_{\diamond}^{2}$ and $\mathbb{Z}_{m}^{2}$. The appropriate
neighborhoods of origin are measured in the norm $\|x + \ii y\|_1 = |x|+|y|$,
and we denote these neighborhoods on the two lattices by
\begin{align*}
\overline{B}_\diamed(r) = \left\{ z \in \ZZ^2_\diamed \, : \, \|z\|_1 \leq r \right\} .
\end{align*}


\begin{lem}\label{lem: discrete Laurent monomials}
There exist unique functions $z\mapsto z^{\left[k\right]}$, $k\in\mathbb{Z}$,
defined for $z \in \mathbb{Z}_{\diamond}^{2} \cup \mathbb{Z}_{m}^{2}$ such that:
\begin{itemize}
\setlength{\itemsep}{0.5em}
\item The functions $z\mapsto z^{\left[k\right]}$ for $k \in \{-1,0,1\}$ are given by
\[z^{\left[0\right]} \equiv 1 , \qquad z^{[1]}=z \]
and
\[ z^{[-1]} = \begin{cases} 2\pi \, \Cker(z) & \text{ for } z \in \ZZ^2_m \\
\frac{\pi}{2} \underset{a \in \{ \pm \frac{1}{2} , \pm \frac{\ii}{2} \}}{\sum} \Cker(z-a)
& \text{ for } z \in \ZZ^2_\diamond \end{cases} .\]
\item For all $k\in\bZ$ we have
\begin{align}\label{dr}
[\dd z^{\left[k\right]} ] = k z^{\left[k-1\right]}.
\end{align}
\item For $k \geq 0$ the function $z \mapsto z^{[k]}$ is discrete holomorphic in the sense
that \[ [\ddbar z^{[k]}] \equiv 0 . \]
\item For all $k>0$ we have the vanishing at origin conditions
\begin{align*}
0^{[k]}=0 , \qquad
\left( \frac{1}{2} \right)^{[k]} + \left( \frac{-1}{2} \right)^{[k]} = 0 , \qquad
\left( \frac{\ii}{2} \right)^{[k]} + \left( \frac{-\ii}{2} \right)^{[k]} = 0 , & \\ %\qquad
\text{ and } \; \left( \frac{1+\ii}{2} \right)^{[k]} + \left( \frac{1-\ii}{2} \right)^{[k]}
+ \left( \frac{-1+\ii}{2} \right)^{[k]} + \left( \frac{-1+\ii}{2} \right)^{[k]} = 0 . &
\end{align*}
%$0^{[k]}=0$, and $(\frac{1}{2})^{[k]} + (\frac{-1}{2})^{[k]} = 0$,
%and $(\frac{\ii}{2})^{[k]} + (\frac{-\ii}{2})^{[k]} = 0$, and
%$(\frac{1+\ii}{2})^{[k]} + (\frac{-1+\ii}{2})^{[k]} + (\frac{-1-\ii}{2})^{[k]} + (\frac{1-\ii}{2})^{[k]} = 0$.
\end{itemize}
The functions $z \mapsto z^{[k]}$ are called the Laurent monomials, and they have the following further properties:
\begin{itemize}
\item For all $k \in \ZZ$ and $z \in \ZZ^2_\diamond \cup \ZZ^2_m$
we have $\overline{z}^{\left[k\right]} = \overline{z^{\left[k\right]}}$.
\item For all $k \in \ZZ$ and $z \in \ZZ^2_\diamond \cup \ZZ^2_m$
we have $(\ii z)^{\left[k\right]} = \ii^k z^{\left[k\right]}$.
\item For $k>0$ the function $z \mapsto z^{[k]}$ vanishes for $z \in \clballdiamed{\frac{k-1}{2}}$.
\item For $k \leq -1$ the function $z \mapsto z^{[k]}$ is
discrete holomorphic outside a neighborhood of the origin in the sense
that
$ [\ddbar z^{[k]}] = 0 $ % \qquad \text{for }
for $ z \notin \clballdiamed{\frac{|k|}{2}} $.
%\item{For $k \ge 2$,
% \[0^{\left[k\right]}=\left(\pm\frac{1}{2}\right)^{\left[k\right]}=\left(\pm\frac{i}{2}\right)^{\left[k\right]}=0.\]}
\end{itemize}
\end{lem}
In the proof below we use discrete integrations that differ slightly from the important integrals later
(for the attentive reader there should be no confusion, however, as later we will use
integration contours on a different lattice, and we will integrate products of two functions).
Below, for a function $f$ defined at the midpoints of edges of $\ZZ^2$ (i.e., on $\ZZ^2_m$),
we define the integral along the oriented edge $\overrightarrow{zw}$ ($z,w \in \ZZ^2$, $|z-w|=1$)
by the evident expression $(w-z) \times f(\frac{z+w}{2})$.
%\[ \int_{\overrightarrow{zw}}^{(\ZZ^2)} f \; = \; (w-z) \times f \left( \frac{z+w}{2} \right) . \]
The integral along an oriented path is the sum of the integrals along the oriented edges of the path.
Note that the positively oriented integral around a plaquette surrounding $z \in \ZZ^2_*$
of $f\colon \ZZ^2_m \rightarrow \bC$ is just $2\ii [\ddbar f](z)$.
The set of horizontal edges $(\ZZ+\frac{1}{2}) \times \ZZ \subset \ZZ^2_m$,
the set of vertical edges $\ZZ \times (\ZZ+\frac{1}{2}) \subset \ZZ^2_m$, and the
dual $\ZZ^2_* = (\ZZ+\frac{1}{2})^2$ are all translates of the standard square grid $\ZZ^2$,
and we define integrals on these with obvious modifications.
\begin{proof}[Proof of Lemma~\ref{lem: discrete Laurent monomials}]
The case of negative $k$ is obvious once $z^{[-1]}$ has been defined by the formula given above.
Indeed, to have the correct derivatives of Equation \ref{dr}, we need to set \[ z^{[-1-m]} = \frac{(-1)^m}{m!} [\dd^m z^{[-1]}] . \]
It follows from the property $[\ddbar z^{[-1]}] = 2\pi \delta_0(z) + \frac{\pi}{2} \sum_a \delta_a(z)$
and the definition of the difference operator $\dd$, that $[\ddbar z^{[-1-m]}]$ vanishes
outside the neighborhood $\clballdiamed{\frac{1+m}{2}}$ of the origin.
The symmetry $(\ii z)^{[-1-m]} = \ii^{-1-m} z^{[-1-m]}$ follows
from the symmetry $\potker(\ii z) = \potker(z)$ of the potential kernel, and the expression for $\dd$.
The symmetry $\overline{z}^{\left[-1-m\right]} = \overline{z^{\left[-1-m\right]}}$ similarly follows
from $\potker(\overline{z}) = \potker(z)$ and the expression for $\dd$.

For positive $k$, we proceed recursively.
We first define $z^{[0]}$ and $z^{[1]}$ by the given formulas (they obviously satisfy the stated properties).
Then for $k \geq 2$ we define $z^{[k]}$ by integrating $k z^{[k-1]}$.
For $z \in \ZZ^2$ the path of integration is any path from $0\in\ZZ^2$ to $z$.
The result is independent of the choice of such path because $[\ddbar z^{[k-1]}] \equiv 0$.
For $z$ (a midpoint of) a horizontal edge the path of integration is any path from
$\frac{1}{2} \in (\ZZ+\frac{1}{2}) \times \ZZ$ to $z$
(we could alternatively take the starting point $-\frac{1}{2}$ without any change, since $0^{[k-1]} = 0$),
and for $z$ (a midpoint of) a vertical edge the path of integration is any path from
$\frac{\ii}{2} \in \ZZ \times (\ZZ+\frac{1}{2})$ to $z$
(we could alternatively take $-\frac{\ii}{2}$ as the starting point).
For $z \in \ZZ^2_*$ the path of integration is any path from
$\frac{1+\ii}{2} \in \ZZ^2_*$ to $z$, and we include an additive constant so as to
guarantee $\sum_{a \in \{\pm\frac{1+\ii}{2},\pm\frac{1-\ii}{2}\}} a^{[k]} = 0$
(in fact, this additive constant is zero except when $k=2$, in which case we have to set
$(\frac{1+\ii}{2})^{[2]}=\frac{\ii}{2}$).

To calculate $[\dd z^{[k]}]$, write it as
\[ [\dd z^{[k]}] =
\frac{1}{2} \Bigg[ \left( z+\frac{1}{2} \right)^{[k]} - \left( z-\frac{1}{2} \right)^{[k]} \Bigg]
- \frac{\ii}{2} \Bigg[ \left( z+\frac{\ii}{2} \right)^{[k]} - \left( z-\frac{\ii}{2} \right)^{[k]} \Bigg] \]
and note that the two terms are both equal to $\frac{1}{2} k z^{[k-1]}$ by the construction
of $z^{[k]}$ as an integral.
Similarly, in the calculation of $[\ddbar z^{[k]}]$ the two contributions cancel.
We have thus constructed Laurent monomials with the desired defining properties,
and uniqueness is clear since only constant functions are annihilated by both $\dd$ and $\ddbar$.

The fact that $z^{[k]}$ is $0$ in $\clballdiamed{\frac{k-1}{2}}$ follows inductively:
the vanishing of $z^{[k-1]}$ in a neighborhood of origin makes its integral $z^{[k]}$ constant in
a slightly bigger neighborhood, and this constant is fixed by the vanishing at origin conditions.
The symmetries $(\ii z)^{[k]} = \ii^k z^{[k]}$ and
$\overline{z}^{\left[k\right]} = \overline{z^{\left[k\right]}}$ are clear by the construction.
\end{proof}
%\begin{rem}\label{rem:finitely-many-modes}
%The following easily verified observation will be important in the Sugawara construction of the $L_n$ operators: for each fixed $z \in \ZZ^2_\diamond \cup \ZZ^2_m$, there exists $N=N(z) < \infty$ such that $z^{[k]} =0$ for all $k \ge N$.
%\end{rem}
%\begin{rem}
%There is some freedom in how to define $z^{[-1]}$ on $\ZZ^2_\diamond$.
%Indeed, we could just as well have defined it as a translate of the Cauchy kernel $\Cker(z)$ on $\ZZ^2_m$,
%thereby placing the singularity at some point of $\ZZ^2_m$. We have chosen instead a symmetric version
%by averaging translates with singularities at $\{\pm 1/2, \pm i/2 \}$.
%\end{rem}
%\begin{rem}
%The vanishing at origin conditions for $k>0$ are only to fix some constants of integration.
%Similarly, there would be some freedom in the definition of $z^{[-1]}$ for $z \in \ZZ^2_\diamond$.
%In both cases, we have chosen to fix the constants in a simple symmetric manner.
%For the results of the present article concerning the discrete GFF in the full plane
%the choices are not important, and concerning the discrete GFF in the half-plane, only
%a symmetry in the reflection $z \mapsto \overline{z}$ would be needed.
%\end{rem}

%Fact
%that for given $z$, only a finite number of $z^{\left[k\right]}$
%for $k>0$ are nonzero -- refer to remark to be made about positive
%energy property{]} {[}Symmetries of $z$: $\overline{z}^{\left[k\right]}=\overline{z^{\left[k\right]}}$,
%$\left(-z\right)^{\left[k\right]}=\left(-1\right)^{k}z^{\left[k\right]}$
%-- hopefully correct{]}

\subsection{Discrete contour integrals of products of two functions}
We now define the important notion of discrete contour integral of a product of two functions.
The contour of integration will be a path on the shifted half-mesh square lattice
$(\frac{1}{2}\ZZ + \frac{1}{4})^2$. Note that any edge of this lattice is between a point of
$\ZZ^2_m$ and a point of $\ZZ^2_\diamond$, both at distance $\frac{1}{4}$ of the edge.
\begin{figure}
\begin{center}
\begin{minipage}{13cm}\epsfig{file=pics-lattice_Virasoro.2, width=13cm}
\end{minipage}
\end{center}
\caption{A discrete contour is
a closed counterclockwise oriented path on the shifted
half-mesh square lattice $(\frac{1}{2}\ZZ^2+\frac{1}{4})^2$.}
\label{fig: discrete contours}
\end{figure}
Let $f \colon \ZZ^2_m \rightarrow \bC$ and $g \colon \ZZ^2_\diamond \rightarrow \bC$ be two functions.
We define the integral of the product of $f$ and $g$ along an oriented edge $[\overrightarrow{uv}]$
($u,v \in (\frac{1}{2}\ZZ + \frac{1}{4})^2$, $|u-v|=\frac{1}{2}$) as
\[ \int_{[\overrightarrow{uv}]} f(z) g(z) \ud z \; := \; (v-u) \times f(z_m) g(z_\diamond) ,\]
where $z_m \in \ZZ^2_m$ and $z_\diamond \in \ZZ^2_\diamond$ are the points at distance $\frac{1}{4}$
of $[\overrightarrow{uv}]$.
The integral along an oriented path $[\gamma]$ on $(\frac{1}{2}\ZZ + \frac{1}{4})^2$
is the sum of the integrals along the oriented edges of the path, and is denoted by
$ \int_{[\gamma]} f(z) g(z) \ud z $
(by convention, we write the function defined on $\ZZ^2_m$ first, and the function defined
on $\ZZ^2_\diamond$ second).

By discrete contour we mean a positively oriented
(i.e. counterclockwise) closed simple path $[\gamma]$ on $(\frac{1}{2}\ZZ + \frac{1}{4})^2$.
We denote by
\[ \mathrm{int}_\diamond[\gamma] \subset \ZZ^2_\diamond \qquad (\text{respectively }
\mathrm{int}_m[\gamma] \subset \ZZ^2_m ) \]
the set of points of the diamond lattice (respectively the medial lattice) encircled by $[\gamma]$.
The discrete closure
$\overline{\mathrm{int}}_\diamond[\gamma] \subset \ZZ^2_\diamond$ (respectively
$\overline{\mathrm{int}}_m[\gamma] \subset \ZZ^2_m$) is the
union of $\mathrm{int}_\diamond[\gamma] \subset \ZZ^2_\diamond$ (respectively
$\mathrm{int}_\diamond[\gamma] \subset \ZZ^2_\diamond$) and
the set of points of $\ZZ^2_\diamond$ (respectively $\ZZ^2_m$)
at distance $\frac{1}{4}$ from $[\gamma]$.
The integrals over such closed discrete contours are denoted by
\[ \oint_{[\gamma]} f(z) g(z) \ud z , \]
and they can be evaluated with the help of the following
combination of discrete Leibnitz rule and Stokes' formula.
\begin{lem}\label{IBP}
Let $f:\ZZ^2_m \to \mathbb{C} $ and $g: \ZZ^2_\diamond \to \mathbb{C}$ and
$[\gamma]$ a discrete contour. Then we have
\[ \oint_{[\gamma]} f(z) g(z) \ud z =
\ii \sum_{z \in \mathrm{int}_m[\gamma]} f(z) \; [\ddbar g](z)
+ \ii \sum_{z \in \mathrm{int}_\diamond[\gamma]} [\ddbar f](z)\; g(z) . \]
%\begin{align*}
% \frac{1}{2i} \oint_{[\gamma]} f(z) \, g(z) dz=\sum_{z \in \cV_{\Gamma_m}} f(z) \left[\ddbar g \right](z) + \sum_{z \in \cV_{\Gamma_\diamond}} \left[\ddbar f\right](z) g(z),
%%\left( \sum_{\vec{e} \in \partial_0 \mathcal{E}_{\Omega}} f(\mathbf{m}(\vec{e})) \, g\left(\mathbf{m}(\vec{e}) + \frac{i \vec{e}}{2}\right)\vec{e} + \sum_{\vec{e} \in  \partial_1 \mathcal{E}_{\Omega^*}} f(\mathbf{m}(\vec{e})) \, g\left(\mathbf{m}(\vec{e}) + \frac{i \vec{e}}{2}\right)\vec{e} \right),
%\end{align*}
%where $\left[ \gamma \right] = \left[ \dd \Gamma \right]$.
\end{lem}
%\begin{rem}
%There are two points per unit area of both $\ZZ^2_m$ and $\ZZ^2_\diamond$,
%so the right hand side is a discretization of the intergal
%$2\ii \iint \bar{\partial} (fg)$ over the area encircled by $\gamma$.
%\end{rem}
\begin{proof}
Decompose the integral to a sum of integrals around elementary plaquettes
of the shifted half-mesh square lattice $(\frac{1}{2}\ZZ + \frac{1}{4})^2$,
and notice that the integral around a plaquette that surrounds $z \in \ZZ^2_m$
is $\ii f(z) \, [\ddbar g](z)$ and the integral around a plaquette that
surrounds $z \in \ZZ^2_\diamond$ is $\ii [\ddbar f](z) \, g(z)$.
\end{proof}



\subsection{Discrete residue calculus}
As a consequence of the above simple lemma,
we have the following discrete analog of a standard residue calculation.
\begin{lem}\label{key-lemma}
Let $m,n \in \ZZ$, let $r = \max \{0, -\frac{m}{2}, -\frac{n}{2} \}$, and let $[\gamma]$ be a discrete contour that separates
the neighborhood $\clballdiamed{r}$ of the origin from infinity.
Then we have
\begin{align}\label{oct31_1}
\frac{1}{2 \pi \ii} \oint_{[\gamma]} z^{[m]} \, z^{[n]} \, \ud z = \delta_{m+n, -1} ,
\end{align}
where $z^{[m]}$ and $z^{[n]}$ are discrete Laurent monomial functions
restricted to $\ZZ^2_m$ and $\ZZ^2_\diamond$, respectively.
%where the function $z^{[m]}\Big|_{\ZZ^2_m}$
%      \begin{align}\label{oct31_1}
%      \frac{1}{2i} \oint_{[\gamma]} z^{[m]} \, z^{[n]} dz = \delta_{m+n+1}.
%      \end{align}
\end{lem}
\begin{proof}
By Lemma~\ref{IBP} we have
\begin{equation}\label{feb26}
\frac{1}{\ii} \oint_{[\gamma]} z^{[m]} \, z^{[n]} \, \ud z
= \sum_{z \in \mathrm{int}_m[\gamma]} z^{[m]} \, [\ddbar z^{[n]} ]
+ \sum_{z \in \mathrm{int}_\diamond[\gamma]} [\ddbar z^{[m]} ] \, z^{[n]} ,
\end{equation}
and we note that by the assumption on $[\gamma]$, the sets $\mathrm{int}_m[\gamma]$ and
$\mathrm{int}_\diamond[\gamma]$ contain the supports of $[\ddbar z^{[n]}]$
and $[\ddbar z^{[m]}]$, respectively.
Also note that if $m,n \ge 0$, then \eqref{oct31_1} follows immediately from \eqref{feb26},
since $z^{[n]}$ and $z^{[m]}$ are discrete holomorphic.

Suppose $n \geq 0, m \leq -1$. Then only the second sum on the right hand side
of \eqref{feb26} remains.
%Moreover, the sum has no non-vanishing terms unless
%the radius of the neighborhood
%where $z^{[n]}$ vanishes is smaller than the radius of the neighborhood where $[\ddbar z^{[m]}]$ is supported
%$n < |m|+1$.
Using $\dd^\top = -\dd$ repeatedly we get that
\begin{align}
\sum_{z \in \mathrm{int}_\diamond[\gamma]} [\ddbar z^{[m]}] \,  z^{[n]}
& = \frac{(-1)^{-m-1}}{(-m-1)!} \sum_{z \in \mathrm{int}_\diamond[\gamma]} [\ddbar \dd^{-m-1} z^{[-1]}] \, z^{[n]} \\
\nonumber
& = \frac{1}{(-m-1)!} \sum_{z} [\ddbar z^{[-1]}] \, [\dd^{-m-1} z^{[n]}] \\
\nonumber
& = \frac{n!}{(-m-1)!(n+m+1)!} \sum_{z} [\ddbar z^{[-1]}] \, z^{[n+m+1]} ,
\end{align}
where the sums in the last two expressions are over $z \in \ZZ^2_\diamond$ or $z \in \ZZ^2_m$,
according to whether $-m-1$ is even or odd, respectively. In the two cases, we have
that $[\ddbar z^{[-1]}]$ is $2\pi \delta_0(z)$ or
$\frac{\pi}{2} \sum_{a \in \{\pm\frac{1}{2}, \pm\frac{\ii}{2}\}} \delta_a(z)$, respectively,
and in either case the result of the whole expression evaluates to $\delta_{n+m+1}$.
The case $n \leq -1 , m \geq 0$ is similar.

Finally, in the case $m,n \leq -1$ rewrite both sums on the right hand side of \eqref{feb26}
as above. If $n$ and $m$ have the same parity, then each of the rewritten sums vanishes
separately because the summands are odd:
$[\ddbar \Cker](-z) \times (-z)^{[n+m+1]} = - [\ddbar \Cker](z) \times z^{[n+m+1]}$.
If $n$ and $m$ are of different parities, then both rewritten sums are sums over the same
sublattice --- they are otherwise identical, but they come with opposite signs and cancel
each other.
\end{proof}
\begin{rem}
Discrete holomorphic monomials and related constructions have been considered by several authors, apparently beginning with the work of Duffin \cite{duffin}. We mention also the more recent work of Mercat \cite{mercat} which was also in part motivated by applications to statistical mechanics models.
\end{rem}


%\red{[Connection with disintegration of measures, universal property
%of conditional expectation? -- not mandatory to include that]}

\section{Discrete current modes}\label{sect:current-modes}

\subsection{The discrete Gaussian free field pinned at the origin}
While the (massless) discrete Gaussian free field can not directly be defined in the
full plane, the massive Green's function $G_\mass$ introduced in Section~\ref{sect:monomials}
is the covariance of the well-defined massive discrete Gaussian free
field $\left( \phi_\mass(z) \right)_{z \in \ZZ^2}$,
\[ \lra{\phi_\mass(z) \, \phi_\mass(w)} = G_\mass(z-w) . \]
In the massless limit $\mass \searrow 0$, differences of values of the field remain well-behaved.
The pinned discrete Gaussian free field $\phi$ is the limit as $\mass \searrow 0$ of the
field $\phi_\mass -\phi_\mass(0)$. We can write its correlation functions in terms of the
potential kernel $\potker$ introduced in Section~\ref{sect:monomials},
\begin{align*}
\lra{ \phi(z) \, \phi(w) }
= \; & \lim_{\mass \searrow 0} \lra{(\phi_\mass(z) - \phi_\mass(0)) \, (\phi_\mass(w) - \phi_\mass(0))} \\
%= \; & \lim_{\mass \searrow 0} \Big( G_\mass(z-w) - G_\mass(z) - G_\mass(w) + G_\mass(0) \Big) \\
= \; & \potker(z) + \potker(w) - \potker(z-w) .
\end{align*}

\subsection{The discrete current}
Let $\phi$ be the discrete GFF on $\ZZ^2$, pinned at $0$. To define the
\emph{current}
\[ J(z) = [\dd \phi](z) \qquad (z \in \ZZ^2_m) , \]
we first extend the definition of $\phi$ to $\ZZ^2_\diamond = \ZZ^2 \cup \ZZ^2_*$ by
setting the value of $\phi$ to zero on $\ZZ^2_*$. By differentiating the two point correlation function
\[ \lra{\phi(z) \phi(w)} = \potker(z) + \potker(w) - \potker(z-w) \]
%(where $\potker$ is the potential kernel, extended to zero on $\ZZ^2_*$ as in Section~\ref{sect:monomials})
we can now conveniently write the correlations of the current in terms of $\Cker = [\dd \potker]$. For example:
%$z^{[-1]} = 2\pi \Cker(z) = 2\pi [\dd \potker](z)$ (for $z \in \ZZ^2_m$):
\begin{align*}
\lra{J(z) \phi(x)} = \; & \begin{cases}
\Cker(z) - \Cker(z-x) & \text{ if } x \in \ZZ^2 , \;  z \in \ZZ^2_m; \\
0 & \text{ if } x \in \ZZ^2_* , \;  z \in \ZZ^2_m \end{cases}; \\
\lra{J (z)J (w) } = \; &
%\begin{cases}
%\frac{1}{2} \Big( \Cker(w+\frac{1}{2}-z) - \Cker(w-\frac{1}{2}-z) \Big) & \text{ if } w \in \ZZ^2_m \text{ is horizontal} \\
%\frac{1}{2\ii} \Big( \Cker(w+\frac{\ii}{2}-z) - \Cker(w-\frac{\ii}{2}-z) \Big) & \text{ if } w \in \ZZ^2_m \text{ is vertical} \end{cases} \\
[\dd \Cker] (w-z) .
%= \; &
%\begin{cases}
%[\dd^\mathrm{hor} \Cker] (w-z) & \text{ if } w \in \ZZ^2_m \text{ is horizontal} \\
%[\dd^\mathrm{ver} \Cker] (w-z) & \text{ if } w \in \ZZ^2_m \text{ is vertical}. \end{cases}
\end{align*}

\subsection{Wick's formula}
Let $\phi$ be the discrete GFF on $\ZZ^2_\diamond$, pinned at $0$. We have given
the two-point correlations for its current $J = [\dd \phi ]$ in terms of
$[\dd \Cker]$. Wick's formula lets us compute higher order correlations of a
Gaussian field from its two-point correlations; see \cite[Chapter 1]{janson} for a proof.

\begin{lem}[Bosonic Wick's formula]
Suppose that $(X_i)_{i \in I}$ is a finite collection of centered Gaussian random variables. Then we have
\begin{align}\label{wick}
\lra{\prod_{i\in I} X_i } = \sum_{P} \prod_{\{a,b\} \in P} \lra{X_{a} X _{b}},
\end{align}
where the sum is over all pair partitions $P$ of $I$, i.e. over
collections $P$ of $\frac{|I|}{2}$ mutually disjoint two element subsets of $I$
(the sum is empty if $|I|$ is odd).
\end{lem}

We will make calculations involving correlations of the discrete
current $J(z)$ and the field $\phi(x)$. In particular, the Wick expansion of
\[ \Biglra{ J (z) J (w) \phi(x_1) \cdots \phi(x_k)} \]
contains terms of two types: those in which the two currents are paired together
\begin{align}
\label{type2}
\biglra{J (z) J(w)} & \prod_{\{a,b\}\in P} \biglra{\phi(x_{a}) \phi(x_{b})} \\
\nonumber
& \text{($P$ a pair partition of $\set{1,2,\ldots,k}$)}
\end{align}
and those in which both are paired with something else
\begin{align}
\label{type1}
\biglra{J (z) \phi(x_i)} & \biglra{J (w) \phi(x_j)} \prod_{\{a,b\}\in P'} \biglra{\phi(x_{a}) \phi(x_{b})} \\
\nonumber
& \text{($i \neq j$ and $P'$ a pair partition of $\set{1,2,\ldots,k}\setminus\set{i,j}$)} .
\end{align}

\subsection{Discrete current modes acting on insertions}
%Recall that for $z,w \in \ZZ^2_m, \, x \in \ZZ^2_\diamond$ we have
%\[ \lra{J (z) \phi(x)} =  z^{[-1]} - (z-x)^{[-1]} , \quad  \lra{J (z)J (w) }=- [\dd_w(z-w)^{[-1]}]. \]
%\begin{defn}\label{current-mode-def}
Suppose that $\mathcal{G} \subset \ZZ^2$ is finite.
For any $\mu \in \mathcal{M}$ and any $n \in \mathbb{Z}$ we define discrete current
modes $\mathfrak{a}_n$ acting by discrete contour integrals on insertions of
functions $f: \mathbb{R}^{\mathcal{G}} \to \mathbb{C}$ in $L^p(\mu_{\rm{GFF}} \big|_{\mathcal{G}})$
for some $p >1$:
\begin{equation}\label{A_def}
\lra{\mathfrak{a}_n f(\phi \big|_{\mathcal{G}}) }_\mu :=
\frac{1}{\sqrt{\pi}} \oint_{[\gamma]} \lra{ J(z) f(\phi \big|_{\mathcal{G}})g(\phi \big|_{\mathcal{G}})}
    z^{[n]} \, \ud z
\end{equation}
where $[\gamma]$ is any discrete contour separating
$\mathcal{G} \cup \clballdiamed{\max\{0,-\frac{n}{2}\}} $
from $\infty$, and $g=g^\mu_{\mathcal{G}}$ is the Radon-Nikodym derivative of
$\mu \big|_{\mathcal{G}}$ with respect to $\mu_{\rm{GFF}}\big|_{\mathcal{G}}$.

%\red{Sorry, I'm not yet completely sure of good normalization, probably $\alpha = \frac{1}{\sqrt{\pi}}$\ldots}

More generally, we define the iterated action (or ``product'') of multiple
discrete current modes on field insertions by contour integrals
\begin{multline*}
\lra{\mathfrak{a}_{n_j}\cdots\mathfrak{a}_{n_1}  f(\phi \big|_{\mathcal{G}}) }_\mu \\
:= \frac{1}{{\pi}^{j/2}} \oint_{\left[\gamma_{j}\right]} \cdots \oint_{\left[\gamma_{1}\right]}
    \lra{ J(w_j) \cdots J(w_1) f(\phi \big|_{\mathcal{G}})g(\phi \big|_{\mathcal{G}})}
    w_1^{\left[n_1\right]} \cdots w_j^{\left[n_j\right]} \ud w_1 \cdots \ud w_j,
\end{multline*}
where the contours are radially ordered: the contour $[\gamma_1]$ of integration of $w_1$
%(corresponding to $\mathfrak{a}_{n_1}$)
separates
$\mathcal{G} \cup \clballdiamed{\max\{0,-\frac{n_1}{2}\}}$ from $\infty$,
and for $k=2,3,\ldots,j$, the contour $[\gamma_k]$ of integration of $w_k$
%(corresponding to $\mathfrak{a}_{n_k}$)
separates $\overline{\mathrm{int}}_m [\gamma_{k-1}] \cup
\overline{\mathrm{int}}_\diamond [\gamma_{k-1}] \cup \clballdiamed{\max\{0,-\frac{n_k}{2}\}}$
%(corresponding to $\mathfrak{a}_{n_{k-1}}$)
from $\infty$.
%\end{defn}

%\begin{rem}
%We are mostly interested in the case when $f$ is a polynomial.
%\end{rem}
%\begin{rem}
%In view of the residue formula of Lemma~\ref{key-lemma} it would perhaps be more accurate to call $\mathfrak{j}_n=\frac{1}{2i}\oint J(z) z^{[n]} dz$ current mode, so that $\mathfrak{a}_n = i\mathfrak{j}_n$.
%\end{rem}



%\begin{prop}
%\label{prop:contour-integral-current-mode}For any $x_{1},\ldots,x_{k}\in\mathbb{Z}_\diamond^{2}$
%and any $n\in\mathbb{Z}$ the discrete contour integral \eqref{A_def}
%is independent of the discrete integration contour $[\gamma]$ provided it separates
%$\left\{0, x_{1},\ldots,x_{k}\right\}$ from $\infty$.
%\end{prop}
%\begin{proof}
%Suppose $[\gamma']$ is a discrete contour that separates $[\gamma]$ from infinity.
%The integration by parts formula of Lemma~\ref{IBP} shows that for any $x \in \ZZ^2_\diamond$ inside $[\gamma]$,
%\[
%\frac{1}{2i}\left(\oint_{\left[\gamma' \right]}  - \oint_{\left[\gamma\right]} \right)(z-x)^{[-1]} z^{\left[n\right]}\mathrm{d}z = \sum_{z \in [\Gamma' \setminus \Gamma]_m} (z-x)^{[-1]} [\ddbar_z z^{\left[n\right]}] + \sum_{z \in [\Gamma' \setminus \Gamma]_\diamond} [\ddbar_z(z-x)^{[-1]}] z^{\left[n\right]} =0
%\]
%by discrete holomorphicity. A Wick expansion combined with this computation completes the proof.
%\end{proof}
%It is no more difficult to check this directly for $L^2$ insertions, so we will do so.x

It follows from H\"older's inequality that the integrals %inside the integrals
$ \lra{\mathfrak{a}_{n} f(\phi |_{\mathcal{G}}) }_\mu $ and
$ \lra{\mathfrak{a}_{n_j}\cdots\mathfrak{a}_{n_1} f(\phi |_{\mathcal{G}}) }_\mu $ make
sense for any given choice of contours, but we need to verify that the definition is
independent of the choices of contours.
\begin{lem}\label{prop:contour-integral-current-mode}
Suppose that $\mathcal{G} \subset \ZZ^2$ is finite and
that $f: \mathbb{R}^{\mathcal{G}} \to \mathbb{C}$ is
in $L^p(\mu_{\rm{GFF}} \big|_{\mathcal{G}})$ for some $p>1$. Then,
\[ \oint_{\left[\gamma_{j}\right]} \cdots \oint_{\left[\gamma_{1}\right]}
  \lra{ J(w_j) \cdots J(w_1) f(\phi \big|_{\mathcal{G}})}
  w_1^{\left[n_1\right]} \cdots w_j^{\left[n_j\right]} \ud w_1 \cdots \ud w_j, \]
is independent of the sequence of contours $[\gamma_1] , [\gamma_2], \ldots, [\gamma_j]$,
provided the contours are radially ordered as above.
%with $[\gamma_1]$ separating $\mathcal{G} \cup \left\{0, \pm 1/2, \pm i/2\right \}$ from $\infty$.
%the discrete contour integral
%\[
%\oint_{[\gamma]}\lra{J(z) f(\phi\big|_{\mathcal{G}})} z^{[n]} dz,
%\]
%is independent of the choice of bounded integration contour $[\gamma]$ provided it separates
%$\mathcal{G} \cup \left\{0, \pm 1/2, \pm i/2\right \}$ from $\infty$.
\end{lem}
%For any finite $\mathcal{G} \subset \ZZ^2_\diamond$, any $n\in\mathbb{Z}$, and any $f: \mathbb{R}^{\mathcal{G}} \to \mathbb{C}$ in $L^2\left(\mu_{\rm{GFF}}\big|_{\mathcal{G}}\right)$,

\begin{proof}
%Since $\mu \in \mathcal{M}$ there exists a polynomial $P$ such that
%\[
%\lra{J(z) F(\phi)}_\mu=\lra{J(z) F(\phi) P(\phi\big|_{\mathcal{H}})}.
%\]
We consider the case of one contour first, and we assume $\mathcal{G} \neq \emptyset$.
By the domain Markov property of $\phi$ under $\mu_{\rm{GFF}}$,
if we condition on the values of $\phi$ on $\mathcal{G}$,
the conditional distribution of $\phi$ outside
of $\mathcal{G}$ is that of a Gaussian free field with Dirichlet zero boundary conditions
plus a discrete harmonic function $h \colon \ZZ^2 \setminus \mathcal{G} \to \mathbb{C}$
with boundary values on $\dd (\ZZ^2 \setminus \mathcal{G})$ given by the values of $\phi\big|_{\mathcal{G}}$.
Consequently, by conditioning inside the expected value, since $f(\phi\big|_{\mathcal{G}})$ is measurable
with respect to the values of $\phi$ on $\mathcal{G}$, we have
\begin{align}\label{mar4.1}
\lra{J(z) f(\phi\big|_{\mathcal{G}})} = \lra{ [\dd h](z) f(\phi\big|_{\mathcal{G}})},
\qquad z \in \ZZ^2_m .
\end{align}
Now suppose $[\gamma']$ is another positively oriented closed discrete contour that also separates
$\mathcal{G} \cup \clballdiamed{\max\{0,-\frac{n}{2}\}}$ from infinity, and which we may for
simplicity assume to encircle at least the same set as $[\gamma]$.
%which in turn lies outside of $\mathcal{G}$.
The identity \eqref{mar4.1} now shows that
\begin{align*}
\left(\oint_{\left[\gamma' \right]}  - \oint_{\left[\gamma\right]} \right)
    \lra{J(z) f(\phi\big|_{\mathcal{G}})} z^{[n]} \ud z
\; & = \left(\oint_{\left[\gamma' \right]}  - \oint_{\left[\gamma\right]} \right)
    \lra{ [\dd h](z) f(\phi\big|_{\mathcal{G}})} z^{[n]} \ud z \\
\; & = \lra{ f(\phi\big|_{\mathcal{G}}) \times \left(\oint_{\left[\gamma' \right]}  - \oint_{\left[\gamma\right]} \right)
    [\dd h](z) z^{[n]} \ud z } % \\ \; &
= 0 .
\end{align*}
Indeed, the last equality follows since $z^{[n]}$ is discrete holomorphic
outside of $\clballdiamed{-\frac{n}{2}}$, and so for any discrete harmonic function $h$ as above, the %integration by parts
formula of Lemma~\ref{IBP} gives
\begin{align*}
& \frac{1}{\ii}\left(\oint_{\left[\gamma' \right]}  - \oint_{\left[\gamma\right]} \right)[\dd h](z) z^{[n]} \ud z \\
= \; & %\Big(
    \sum_{z \in \mathrm{int}_m[\gamma'] \setminus \mathrm{int}_m[\gamma]}
    %- \sum_{z \in \mathrm{int}_m[\gamma] \setminus \mathrm{int}_m[\gamma']} \Big)
    [\dd h](z) \, \underbrace{[\ddbar z^{\left[n\right]}]}_{=0} %\\
%& \qquad + %\Big(
    + \sum_{z \in \mathrm{int}_\diamond[\gamma'] \setminus \mathrm{int}_\diamond[\gamma]}
    %- \sum_{z \in \mathrm{int}_\diamond[\gamma] \setminus \mathrm{int}_\diamond[\gamma']} \Big)
    \underbrace{[\Delta h](z)}_{=0} \, z^{\left[n\right]} .
%+ \sum_{z \in [\Gamma' \setminus \Gamma]_\diamond} [\Delta u](z) z^{\left[n\right]} = 0.
\end{align*}

It remains to check the case with several radially ordered contours.
By conditioning on the values of $\phi$ on $\mathcal{G}$ and on the
contours $[\gamma_{k-1}]$ and $[\gamma_{k+1}]$ (more precisely, on the vertices of $\ZZ^2$ at
distance $\frac{1}{4}$ from these contours), a similar argument as the one just given implies
that we can move $[\gamma_k]$ as long as it is between $[\gamma_{k-1}]$ and $[\gamma_{k+1}]$.
\end{proof}




%By the domain Markov property of $\phi$ under $\mu_{\rm{GFF}}$, if we condition on the values of $\phi$ on $\mathcal{G}$ (which we can assume is simply connected and contains $\{0, \pm 1/2, \pm i/2\}$), the conditional distribution of $\phi$ outside of $\mathcal{G}$ is that of a mean zero Gaussian free field plus a discrete harmonic function $u(z): \ZZ^2_\diamond \to \mathbb{C}$ with boundary values on $\dd \ZZ^2_\diamond \setminus \mathcal{H}$ given by the values of $\phi$. Hence by conditioning inside the integral, since $F(\phi)$ and $P(\phi\big|_{\mathcal{H}})$ are measurable with respect to the values of $\phi$ on $\mathcal{H}$,
%\begin{align}\label{mar4.1}
%\lra{J(z) F(\phi) P(\phi\big|_{\mathcal{H}})} = \lra{ [\dd u](z) F(\phi) P(\phi\big|_{\mathcal{H}})}.
%\end{align}
%Now suppose $[\gamma']$ is a discrete contour that separates $[\gamma]$ from infinity which in turn is outside of $\mathcal{H}$. The identity \eqref{mar4.1} combined with the integration by parts formula of Lemma~\ref{IBP} shows that
%\begin{align*}
%\left(\oint_{\left[\gamma' \right]}  - \oint_{\left[\gamma\right]} \right)\lra{J(z) F(\phi)}_{\mu} z^{[n]} dz & = \left(\oint_{\left[\gamma' \right]}  - \oint_{\left[\gamma\right]} \right) \lra{ [\dd u](z) F(\phi) P(\phi\big|_{\mathcal{H}})}z^{[n]} dz \\
%&  \lra{ \left(\oint_{\left[\gamma' \right]}  - \oint_{\left[\gamma\right]} \right)[\dd u](z) z^{[n]} dz  F(\phi) P(\phi\big|_{\mathcal{H}})}\\
%&=0.
%\end{align*}
%Indeed, since $z^{[n]}$ is discrete holomorphic outside of $\mathcal{H}$, for any discrete harmonic function $u$ as above,
%\[
%\frac{1}{2i}\left(\oint_{\left[\gamma' \right]}  - \oint_{\left[\gamma\right]} \right)[\dd u](z) z^{[n]} dz = \sum_{z \in [\Gamma' \setminus \Gamma]_m} [\dd u](z) [\ddbar z^{\left[n\right]}] + \sum_{z \in [\Gamma' \setminus \Gamma]_\diamond} [\Delta u](z) z^{\left[n\right]}=0.
%\]
%This completes the proof.
%\end{proof}



\subsection{Current modes as change of measure operators}
Each discrete current mode $\mathfrak{a}_{n}$ induces a change of measure operator acting on $\mathcal{M}$,
denoted by $a_{-n}$.
%Recall that we always assume the corresponding discrete contours to be radially ordered.
%
%Let $\mu \in \mathcal{M}$ and let $P$ and $\mathcal{H}$ be as in Definition~\ref{def:space-gibbs-measures}. Using Lemma~\ref{prop:contour-integral-current-mode} we can consider the contour integral action of the current modes on insertions of $L^2\left(\mu_{\rm{GFF}} \big|_{\mathcal{G}} \right)$ functions $F: \mathbb{R}^{\mathcal{G}} \to \mathbb{C}$ of the free field under the measure $\mu$:
%\[
%\frac{1}{2}\oint_{[\gamma]}\lra{J(z) F(\phi \big|_{\mathcal{G}}) }_\mu z^{[n]} dz=\frac{1}{2}\oint_{[\gamma]}\lra{J(z) F(\phi \big|_{\mathcal{G}} ) P(\phi \big|_{\mathcal{H}}) } z^{[n]} dz,
%\]
%where $[\gamma]$ is any discrete contour that separates $\mathcal{H}$ from $\infty$.

We first state an auxiliary result on the reconstruction of the change of measure operator
from a consistent family of operators acting on insertions. Below we denote briefly
$L^p_\mathcal{G} = L^p(\mu_\mathrm{GFF}\big|_\mathcal{G})$. Note that a Gibbs
measure $\mu \in \mathcal{M}$ is determined by the collection of Radon-Nikodym derivatives
$\left( g_\mathcal{G} \right)_{\mathcal{G} \subset \ZZ^2 \text{ finite}}$,
\[ g_\mathcal{G} = \frac{\ud \mu \big|_\mathcal{G}}{\ud \mu_\mathrm{GFF}\big|_\mathcal{G}}
\in \bigcap_{p<\infty} L^p_\mathcal{G} , \]
which satisfy the consistency conditions that for any $\mathcal{G} \subset \mathcal{G}'$ we have
\[ g_\mathcal{G} (\phi|_\mathcal{G}) =
\EX \left[ g_\mathcal{G'} (\phi|_\mathcal{G'}) \, \Big| \, \phi|_\mathcal{G} \right] ,\]
where the right hand side is the conditional expected value over the discrete GFF $\phi$
given its values on $\mathcal{G}$. Conversely, any collection
$\left( g_\mathcal{G} \right)_{\mathcal{G} \subset \ZZ^2 \text{ finite}}$ of functions
$g_\mathcal{G} \in \bigcap_{p<\infty} L^p_\mathcal{G}$ satisfying the above consistency
condition determines a Gibbs measure $\mu \in \mathcal{M}$.
\begin{lem}\label{lem: lifting to change of measure operators}
Suppose that $\left( \mathfrak{A}_\mathcal{G} \right)_{\mathcal{G} \subset \ZZ^2 \text{ finite}}$
is a family of linear functionals
\[ \mathfrak{A}_\mathcal{G} \colon \bigcup_{p>1} L^p_\mathcal{G} \rightarrow \bC \]
such that
\begin{itemize}
\item for any $p>1$, the functional $\mathfrak{A}_\mathcal{G}$
restricted to $L^p_\mathcal{G}$ is a bounded operator
\item for all $\mathcal{G} \subset \mathcal{G}'$ and $f \in L^p_\mathcal{G}$, the consistency
$ \mathfrak{A}_\mathcal{G'} (f) = \mathfrak{A}_\mathcal{G} (f) $ holds.
\end{itemize}
Then there exists a unique change of measure operator $A \colon \mathcal{M} \rightarrow \mathcal{M}$
such that for any $\mu \in \mathcal{M}$ and $\mathcal{G}$
%with  $\frac{\ud \mu \big|_\mathcal{G}}{\ud \mu_\mathrm{GFF}\big|_\mathcal{G}} = g_\mathcal{G}$
and any $f \in L^p_\mathcal{G}$, $p>1$, we have
\[ \lra{ f }_{A\mu}
= \mathfrak{A}_\mathcal{G} (f g_\mathcal{G}^\mu) .\]
\end{lem}
\begin{proof}
Fix a Gibbs measure $\mu \in \mathcal{M}$. Uniqueness of the Gibbs measure $A\mu$ with the above
correlation functions is clear, so it suffices to construct one and prove
that $A\mu \in \mathcal{M}$.

Fixing first also $\mathcal{G}$, we note that the map
\[ f \mapsto \mathfrak{A}_\mathcal{G} (f g_\mathcal{G}^\mu) \]
is a bounded linear functional on any $L^p_\mathcal{G}$, $p>1$.
It follows from $L^p$-space duality, that there is a unique function
$\alpha^\mu_\mathcal{G} \in \bigcap_{p<\infty} L^p_\mathcal{G}$ such that we can write
\[ \mathfrak{A}_\mathcal{G}(f g_\mathcal{G}^\mu) = \lra{ f \alpha^\mu_\mathcal{G} } .\]
Consistency of the operators implies that for all $f \in L^p_\mathcal{G}$
and for any $\mathcal{G}' \supset \mathcal{G}$ we have
\[ \lra{f \alpha^\mu_{\mathcal{G}'}} = \lra{ f \alpha^\mu_\mathcal{G} } .\]
Thus the collection
$\left( \alpha_\mathcal{G} \right)_{\mathcal{G} \subset \ZZ^2 \text{ finite}}$
is consistent and can be used to define the Gibbs measure $A \mu$ such that the Radon-Nikodym
derivative of $(A \mu)\big|_\mathcal{G}$ is $\alpha_\mathcal{G}$. By construction
$A \mu$ has the desired correlation functions, the Radon-Nikodym derivative
$\alpha^\mu_\mathcal{G}$ is in $L^p_\mathcal{G}$ for all $p<\infty$, and
$(A \mu)\big|_\mathcal{G}$ is absolutely continuous with respect to $\mu\big|_\mathcal{G}$ ---
so we have that $A \mu \in \mathcal{M}$ is a change of measure of $\mu$.
\end{proof}



\begin{prop}\label{prop:current-mode-operators}
For each $n\in\mathbb{Z}$, there
exists a unique operator $a_{n} \colon \mathcal{M}\to\mathcal{M}$
such that for any $\mu\in\mathcal{M}$ and any $x_1, \ldots, x_k \in \mathbb{Z}^{2}$,
\[ \left\langle \phi(x_1) \cdots \phi(x_k)\right\rangle _{a_{n}\mu}= \lra{\mathfrak{a}_{-n} \phi(x_1) \cdots \phi(x_k)}_{\mu}.
%\frac{1}{2}\oint_{\left[\gamma\right]}\left\langle J\left(z\right)\phi(x_1) \cdots \phi(x_k)\right\rangle _{\mu}z^{\left[n\right]}\mathrm{d}z,
\]
The operator $a_n$ is parity reversing.
% odd in the sense that $a_{n}\mathcal{S}\subset\mathcal{A}$ and $a_{n}\mathcal{A}\subset\mathcal{S}$.
Moreover, for any $n_1,\ldots,n_j \in \ZZ$, the composition of operators
$a_{n_1} \cdots a_{n_j} \colon \mathcal{M} \to \mathcal{M}$ is characterized by the formulas
\[ \left\langle \phi(x_1) \cdots \phi(x_k)\right\rangle _{a_{n_1} \cdots a_{n_j}\mu}
= \lra{\mathfrak{a}_{-n_j} \cdots \mathfrak{a}_{-n_1} \phi(x_1) \cdots \phi(x_k)}_{\mu} , \]
for all $\mu\in\mathcal{M}$ and $x_1, \ldots, x_k \in \mathbb{Z}^{2}$, and we recall that the product of
current modes $\mathfrak{a}_{-n_j} \cdots \mathfrak{a}_{-n_1}$ acting on insertions is taken radially ordered.
\end{prop}
\begin{proof}
Since polynomials in $\left( \phi(x) \right)_{x \in\mathcal{G}}$ are dense in any $L^p_\mathcal{G}$,
and the association $f \mapsto \lra{\mathfrak{a}_{-n} f}$ is bounded in any $L^p_\mathcal{G}$,
the existence and uniqueness of the change of measure operator $a_n$ follows from Lemma~\ref{lem: lifting to change of measure operators}.
The formula for the correlation functions of $a_{n_1} \cdots a_{n_j} \mu$ is obtained iteratively, and the
uniqueness of the measure with these correlation functions is again clear.
To show that $a_n$ takes symmetric measures to antisymmetric measures and vice versa, just note
that the Radon-Nikodym derivatives $\alpha^\mu_\mathcal{G}$ for $a_n \mu$, as constructed in Lemma~\ref{lem: lifting to change of measure operators},
have the opposite parity compared to the Radon-Nikodym derivatives $g^\mu_\mathcal{G}$ for $\mu$.
%Let $\mathcal{G} \subset \ZZ^2$ be finite and $\mu \in \mathcal{M}$. Then by definition $\mu\big|_{\mathcal{G}} \ll \mu_{\rm{GFF}}\big|_{\mathcal{G}}$ and $\ud \mu\big|_{\mathcal{G}} = g \, \ud \mu_{\rm{GFF}}\big|_{\mathcal{G}}$, where $g \in L^p=L^p(\mu_{\rm{GFF}}\big|_{\mathcal{G}})$ for every $p < \infty$. Consequently, the assignment
%\begin{equation}\label{def:assignment}
%f \mapsto \lra{\mathfrak{a}_{-n_j} \cdots \mathfrak{a}_{-n_1}  f(\phi \big|_{\mathcal{G}}) g(\phi \big|_{\mathcal{G}}) }
%\end{equation}
%defines for each $q >1$ a bounded linear functional on $L^q$ and by Lemma~\ref{prop:contour-integral-current-mode} it is independent on the choices of contours, provided they are radially ordered. It follows by $L^p$-space duality that there is a unique function $g_{n_1,\ldots,n_j} \in L^p$, for every $p<\infty$, such that
%\begin{align}\label{mar5.1}
%\lra{\mathfrak{a}_{-n_j} \cdots \mathfrak{a}_{-n_1}  f(\phi \big|_{\mathcal{G}}) g(\phi \big|_{\mathcal{G}}) } = \lra{f(\phi \big|_{\mathcal{G}}) g_{n_1,\ldots,n_j}(\phi \big|_{\mathcal{G}}) }
%\end{align}
%for any $f \in L^q$ with $q > 1$.
%%(In fact, since clearly $g=2^{-1}\mathbb{E}\left[ \oint J(z) z^{[n]}Q \mid \phi\big|_{\mathcal{G}}\right]$, we see that $g \in L^p$ for all $p < \infty$.)
%We define a complex Borel measure $(a_{n_1} \cdots a_{n_j}\mu)\big|_{\mathcal{G}}:= g_{n_1,\ldots,n_j} \cdot \mu_{\rm{GFF}}\big|_{\mathcal{G}}$, with $g_{n_1,\ldots,n_j}$ as above.
%Varying $\mathcal{G}$, using the independence on the choices of contours, this defines a complex Gibbs measure $(a_{n_1} \cdots a_{n_j} \mu)$: if $\mathcal{G} \subset \mathcal{G}'$ with densities $g_{n_1,\ldots,n_j}$ and $g_{n_1,\ldots,n_j}'$, then $g_{n_1,\ldots,n_j} = \int_{\phi\big|_{\mathcal{G}' \setminus \mathcal{G}} }g_{n_1,\ldots,n_j}'  \mu_{\rm{GFF}}(\ud \phi)$.
%
%It remains to show that  $a_n$ takes symmetric measures to antisymmetric measures, and vice versa.
%But this follows at once since $J(z)=[\dd \phi](z)$ is a linear combination of linear terms and so
%the density $g_n$ defining $(a_n \mu ) \big|_\mathcal{G}$ has the opposite parity compared to $g$.
\end{proof}
\begin{thm}\label{thm:heisenberg}
The change of measure operators $(a_n), \, n \in \ZZ,$ form a representation of the Heisenberg algebra:
\[ [a_m, a_n] = m \delta_{m+n,0} \, \id_\mathcal{M}. \]
\end{thm}
\begin{proof}
Let $x_1, \ldots, x_k$ be points in $\ZZ^2$ and let $[\gamma_0], [\gamma_1], [\gamma_2]$ be discrete contours such that $[\gamma_0]$ encircles $\{x_1, \ldots, x_k\} \cup \clballdiamed{\max\{0,-\frac{m}{2},-\frac{n}{2}\}}$,
and $[\gamma_1]$ encircles $[\gamma_0]$, and $[\gamma_2]$ encircles $[\gamma_1]$.
%and such that for $j,k =0,1,2$ and $k > j$, $[\gamma_k]$ separates $[\gamma_j]$ from infinity.
We shall start by computing
\begin{multline}\label{diff}
\lra{\mathfrak{a}_m \mathfrak{a}_n \phi(x_1) \cdots \phi(x_k)} - \lra{\mathfrak{a}_n \mathfrak{a}_m \phi(x_1) \cdots \phi(x_k)} \\
= \frac{1}{\pi} \oint_{[\gamma_2]} \oint_{[\gamma_1]}\lra{ J (z) J (w) \phi(x_1) \cdots \phi(x_k)} w^{[n]} z^{[m]} \ud w \ud z  \\
- \frac{1}{\pi} \oint_{[\gamma_1]} \oint_{[\gamma_0]}\lra{ J (w) J (z) \phi(x_1) \cdots \phi(x_k)} z^{[m]} w^{[n]} \ud z \ud w.
\end{multline}
A Wick expansion of the correlation function produces terms of the form \eqref{type1} that factor and thus cancel
in \eqref{diff}, and one non-cancelling term of the form \eqref{type2}. By fixing $w \in \ZZ^2_m$
at distance $\frac{1}{4}$ from $[\gamma_1]$ in each of the two
integrals contributing to \eqref{diff}, and rearranging the integrations with the help of Lemma~\ref{IBP},
we see that the contribution of the non-cancelling terms to \eqref{diff} can be written as
\[ \lra{ \phi(x_1) \cdots \phi(x_k)} \times
\frac{1}{\pi} \oint_{[\gamma_1]} \Big( \oint_{[\sigma_w]} \lra{ J (z) J (w)} z^{[m]} \ud z \Big) w^{[n]} \ud w , \]
where $[\sigma_w]$ is a positively oriented discrete closed contour encircling $w$.
For each $w$, the inner integral therefore becomes
%, according to whether $w$ is the midpoint of a horizontal or vertical edge,
\begin{align*}
\oint_{[\sigma_w]} \lra{ J (z) J (w)} z^{[m]} \ud z
%= \; & \oint_{[\sigma_w]} [\dd^{\mathrm{hor/ver}} \Cker](z-w) \, z^{[m]} \ud z \\
= \; & \oint_{[\sigma_w]} [\dd \Cker](z-w) \, z^{[m]} \ud z \\
%= \; & \ii \sum_{z \in \mathrm{int}_\diamond [\sigma_w]} [\ddbar \dd^{\mathrm{hor/ver}} \Cker](z-w) \, z^{[m]} \\
= \; & \ii \sum_{z \in \mathrm{int}_\diamond [\sigma_w]} [\ddbar \dd \Cker](z-w) \, z^{[m]} \\
%= \; & - \ii \sum_{z \in \overline{\mathrm{int}}_m [\sigma_w]} [\ddbar \Cker](z-w) \, [\dd^{\mathrm{hor/ver}} z^{[m]}] \\
= \; & - \ii \sum_{z \in \overline{\mathrm{int}}_m [\sigma_w]} [\ddbar \Cker](z-w) \, [\dd z^{[m]}] \\
= \; & - \ii \frac{m}{2} w^{[m-1]} ,
\end{align*}
where we used Lemma~\ref{IBP}, the fact that $[\ddbar z^{[m]}] = 0$ inside $[\sigma_w]$ and
the antisymmetry of $\dd$ stated in Lemma~\ref{lem: transpose of dee}. %  = - \dd^{\mathrm{hor/ver}},
%and the observation that
%\[ [\dd^{\mathrm{hor/ver}} w^{[m]}] = \frac{1}{2} [\dd w^{[m]}] = \frac{m}{2} w^{[m-1]} ,\]
%which follows since $[\ddbar w^{[m]}]$ vanishes as $w \notin \clballdiamed{-\frac{m}{2}}$.
We then integrate around $[\gamma_1]$ and apply Lemma~\ref{key-lemma}:
\[ \oint_{[\gamma_1]} \Big( \oint_{[\sigma_w]} \lra{ J (z) J (w)} z^{[m]} \ud z \Big) w^{[n]} \ud w
= \frac{- \ii m}{2} \oint_{[\gamma_1]} w^{[m-1]} w^{[n]} \ud w
= \pi m \delta_{m+n, 0} . \]

Now, let $\mu \in \mathcal{M}$ and a finite $\mathcal{G} \subset \ZZ^2$ be given. The computations above describe the action of $(\mathfrak{a}_m\mathfrak{a}_n -\mathfrak{a}_n\mathfrak{a}_m)$ on insertions of polynomials $P$ in the Gaussian variables
$\left( \phi(x) \right)_{x \in \mathcal{G}}$:
\begin{align}\label{poly}
\lra{(\mathfrak{a}_m\mathfrak{a}_n -\mathfrak{a}_n\mathfrak{a}_m) P(\phi\big|_{\mathcal{G}})}
 % = \; & \pi \alpha^2 m \delta_{0,m+n}\lra{P(\phi\big|_{\mathcal{G}})} \\
 = \; & m \delta_{0,m+n}\lra{P(\phi\big|_{\mathcal{G}})} .
\end{align}
%Let $g=g_\mathcal{G}^\mu$ be the Radon-Nikodym derivative of $\mu\big|_{\mathcal{G}}$ with respect to $\mu_{\rm{GFF}}\big|_{\mathcal{G}}$ and $f \in L^q=L^q\left(\mu_{\rm{GFF}}\big|_{\mathcal{G}}\right)$ for some fixed but arbitrary $q> 1$. H\"older's inequality implies that the product $f g \in L^{p}$ for any $1 \le p <q$. We have noted that polynomials are dense in $L^p(\mu_{\rm{GFF}} \big|_{\mathcal{G}})$ for every $0 \le p < \infty$. Thus whenever $p<q$ we can approximate $fg$ by polynomials in $L^{p}$, and apply H\"older's inequality and \eqref{poly} to see that
%\[ \lra{(\mathfrak{a}_m\mathfrak{a}_n -\mathfrak{a}_n\mathfrak{a}_m) f(\phi \big|_{\mathcal{G}}) g(\phi \big|_{\mathcal{G}}) }
%= m \delta_{0,m+n}\lra{f(\phi \big|_{\mathcal{G}}) g(\phi \big|_{\mathcal{G}})}. \]
%Since $q>1$ and $f \in L^q$ were arbitrary,
Since this holds for all polynomials $P$, we obtain
\[ [a_{-n} , a_{-m}] \, \mu \; = \; m \delta_{0,m+n} \, \mu , \]
and so the proof is complete.
\end{proof}
%\begin{rem}\label{rem:associativity}
%It is not hard to see by similar arguments as those of the proof of Theorem~\ref{thm:heisenberg} that $\mathfrak{a}_k[\mathfrak{a}_m, \mathfrak{a}_n]$ lifts to the change of measure operator $a_k m \delta_{0,m+n} \rm{Id}$.
%\end{rem}
By analogous arguments we may construct antianalytic current modes $\bar{a}_n$
corresponding to $\oint \bar{J}(z) \bar{z}^{[n]} d \bar{z}$, where
$\bar{J}(z):=[\ddbar \phi](z)$, and show that these operators also yield a
representation of the Heisenberg algebra.
%Below we denote by $\overline{\mu}$ the
%complex conjugate of a Gibbs measure $\mu$, i.e. the Gibbs measure defined by
%\[ \lra{P(\phi\big|_\mathcal{G})}_{\overline{\mu}} = \overline{\lra{P(\phi\big|_\mathcal{G})}_{\mu}} \]
%for any polynomial $P$ with real coefficients.
\begin{prop}\label{prop:anti-analytic-sector-an}
For each $k\in\mathbb{Z}$, there
exists a unique operator $\bar{a}_{k}:\mathcal{M}\to\mathcal{M}$, with
$\bar{a}_{k}\mathcal{S}\subset\mathcal{A}$ and $\bar{a}_{k}\mathcal{A}\subset\mathcal{S}$
such that for any $\mu\in\mathcal{M}$ and any (not necessarily distinct)
$x_{1},\ldots,x_{m}\in\mathbb{Z}^{2}_\diamond$ we have
\[ \lra{\phi(x_1) \cdots \phi(x_m)}_{\bar{a}_k\mu} =
    \overline{\lra{\phi(x_1) \cdots \phi(x_m)}_{{a_k} \overline{\mu}}}. \]
Moreover, the change of measure operators $(\bar{a}_n)_{n \in \ZZ}$ form a representation of the Heisenberg algebra:
\[ [\bar{a}_m, \bar{a}_n] = m \delta_{m+n,0} \, \id_\mathcal{M} \]
and they commute with the discrete analytic current modes
\[ [a_m, \bar{a}_n] = 0 . \]
\end{prop}
\begin{proof}
The construction of $\bar{a}_n$ and the calculation of their commutation relations is
identical to those of $a_n$. To see that the two representations of Heisenberg algebra
commute, note that in a calculation like in the proof of Proposition~\ref{thm:heisenberg},
a Wick expansion of $\lra{J(z) \bar{J}(w) \phi(x_1) \ldots \phi(x_k)}$
yields cancelling terms plus a term containing the factor
$\lra{J(z) \bar{J}(w)} = [\ddbar \Cker] (w-z)$. But this function
of $z$ is identically $0$ on the contour $\sigma_w$.
\end{proof}

%\begin{rem}
%We end the section with the observation that the $a_n$ and $\bar{a}_n$ operators commute. Indeed, a Wick expansion of $\lra{J(z) \bar{J}(w) \phi(x_1) \ldots \phi(x_k)}$ yields commuting terms plus a term containing the factor $\lra{J(z) \bar{J}(w)} = -\ddbar_w(z-w)^{[-1]}$. But this function is identically $0$ on the inner integration contour.
%\end{rem}
\section{Sugawara construction of Virasoro modes}\label{sugawara}
We are now ready to construct the operators that span the commuting Virasoro algebras.
\subsection{Virasoro modes acting on insertions}
For any $n \in \ZZ$, the action on insertions of the formally infinite
linear combinations
\begin{align}
\mathfrak{L}_n
& := \frac{1}{2} \sum_{j \ge 0} \mathfrak{a}_{n-j}\mathfrak{a}_j
     + \frac{1}{2}\sum_{j \le -1} \mathfrak{a}_j \mathfrak{a}_{n-j}  \label{Ln-definition} \\
\bar{\mathfrak{L}}_n
& := \frac{1}{2} \sum_{j \ge 0} \bar{\mathfrak{a}}_{n-j}\bar{\mathfrak{a}}_j
     + \frac{1}{2}\sum_{j \le -1} \bar{\mathfrak{a}}_j \bar{\mathfrak{a}}_{n-j}\label{Ln-bar-definition}
\end{align}
of current modes are well defined, since only finitely many terms ever contribute
when we act on insertions of cylinder functions.
Indeed, for fixed finite subgraph $\mathcal{G} \subset \bZ^2$, denote
\[ M_\mathcal{G} = 1 + 2 \times \max \{ \|z\|_1 \; : \; z \in \mathcal{G} \} .\]
Then if $x_1, \ldots, x_k \in \mathcal{G}$, we have
that $\lra{\mathfrak{a}_j \phi(x_1) \cdots \phi(x_k)}$ vanishes for all
$j > M_\mathcal{G}$, because the integration contour
in the definition of $\mathfrak{a}_j$ can then be taken so that in the integration the
monomial $z^{[j]}$ only ever gets evaluated in the set
$\clballdiamed{\frac{j-1}{2}}$, where it is zero. Thus for $f \in L^p_\mathcal{G}$
the action of $\mathfrak{L}_n$ is actually defined as the finite sum
\begin{align*}
\lra{\mathfrak{L}_n f(\phi\big|_\mathcal{G})}
= \; & \frac{1}{2} \sum_{j = 0}^{M_\mathcal{G}} \lra{\mathfrak{a}_{n-j}\mathfrak{a}_j f(\phi\big|_\mathcal{G})}
    + \frac{1}{2} \sum_{j = n - M_\mathcal{G} }^{-1} \lra{\mathfrak{a}_j \mathfrak{a}_{n-j} f(\phi\big|_\mathcal{G})},
\end{align*}
and similarly for $\bar{\mathfrak{L}}_n$.
%with a number of needed terms showing for a fixed $n \in \bZ$ and fixed subgraph $\mathcal{G}$ .

%\begin{rem}
%By Theorem~\ref{thm:heisenberg}, \eqref{Ln-definition} one could equivalently define
%$\mathfrak{L}_n = \frac{1}{2}\mathfrak{a}_n \mathfrak{a}_0
%+ \sum_{j\ge 1} \mathfrak{a}_{n-j}\mathfrak{a}_j$, and similarly for \eqref{Ln-bar-definition}.
%\end{rem}


\subsection{Virasoro modes as change of measure operators}
Like in Section~\ref{sect:current-modes} with current modes, we can lift the action of
$\mathfrak{L}_n$ on insertions to an action on Gibbs measures.
\begin{prop}\label{prop:Virasoro-mode-operators}
For each $n\in\mathbb{Z}$, there
exists unique operators $L_{n}$ and $\bar{L}_n$ on $\mathcal{M}$
such that for any $\mu\in\mathcal{M}$ and any $x_1, \ldots, x_k \in \mathbb{Z}^{2}$
\begin{align*}
\lra{\phi(x_1) \cdots \phi(x_k)}_{L_n \mu}
= \; & \lra{ \mathfrak{L}_{-n} \phi(x_1) \cdots \phi(x_k)}_{\mu} \\
\lra{\phi(x_1) \cdots \phi(x_k)}_{\bar{L}_n \mu}
= \; & \lra{ \bar{\mathfrak{L}}_{-n} \phi(x_1) \cdots \phi(x_k)}_{\mu} .
\end{align*}
The operators $L_n, \bar{L}_n$ are parity preserving.
%even in the sense that $L_{n}\mathcal{S}\subset\mathcal{S}$
%and $L_{n}\mathcal{A}\subset\mathcal{A}$ and similarly for $\bar{L}_n$.
\end{prop}
\begin{proof}
For any $n \in \ZZ$, the operators $\mathfrak{L}_{-n}$ and $\bar{\mathfrak{L}}_{-n}$
form consistent families of operators that are bounded on each $L^p_\mathcal{G}$, $p>1$, as follows
immediately from the fact that for a fixed $\mathcal{G}$ and $n\in\bZ$ they coincide
with finite linear combinations of products of current modes.
Thus by an earlier Lemma, they lift to change of measure operators
$L_{n} \colon \mathcal{M} \rightarrow \mathcal{M}$ and
$\bar{L}_{n} \colon \mathcal{M} \rightarrow \mathcal{M}$, respectively, uniquely determined by
the above formulas on the dense subspace of polynomials.
\end{proof}

Keeping in mind that $a_n$ is the change of measure operator obtained by lifting the
action on insertions of $\mathfrak{a}_{-n}$, and that the lift of a composition
is a composition of lifts in the reverse order (the change of measure operator is a
formal adjoint of the action on insertions), we have a formal expression for the $L_n$ as
\begin{align*}
L_n & = \frac{1}{2} \sum_{j \ge 0} a_{-j} a_{j+n}
     + \frac{1}{2}\sum_{j \le -1} a_{j+n} a_{-j}
   = \frac{1}{2} \sum_{j \le -1} a_{j} a_{n-j}
     + \frac{1}{2}\sum_{j \ge 0} a_{n-j} a_{j} , % \\
%\bar{L}_n & = \frac{1}{2} \sum_{j \ge 0} \bar{a}_{-j} \bar{a}_{j+n}
%     + \frac{1}{2}\sum_{j \le -1} \bar{a}_{j+n} \bar{a}_{-j}
%   = \frac{1}{2} \sum_{j \le -1} \bar{a}_{j} \bar{a}_{n-j}
%     + \frac{1}{2}\sum_{j \ge 0} \bar{a}_{n-j} \bar{a}_{j} .
\end{align*}
and similarly for $\bar{L}_n$.
However, one would have to make sense of these infinite linear
combinations of (compositions of) change of measure operators.
Our approach of lifting the action could be %alternatively
thought of as considering the pointwise convergence of the above series of
change of measure operators when the underlying space of Gibbs measures
is equipped with an appropriate weak topology.
%the (weak) topology induced by integrations against
%finitely supported functions in all $L^p(\mu_{\mathrm{GFF}})$.

%\begin{defn}
%For $n\in \ZZ$, we define the change of measure operators $L_n,\bar{L}_n:\mathcal{M} \to \mathcal{M}$, as the lifts to $\mathcal{M}$ of
%\begin{align}
%\mathfrak{L}_n &:= \frac{1}{2} \sum_{j \ge 0} \mathfrak{a}_{n-j}\mathfrak{a}_j + \frac{1}{2}\sum_{j \le -1} \mathfrak{a}_j \mathfrak{a}_{n-j}; \label{Ln-definition} \\
%\bar{\mathfrak{L}}_n &:= \frac{1}{2} \sum_{j \ge 0} \bar{\mathfrak{a}}_{n-j}\bar{\mathfrak{a}}_j + \frac{1}{2}\sum_{j \le -1} \bar{\mathfrak{a}}_j \bar{\mathfrak{a}}_{n-j}\label{Ln-bar-definition},
%\end{align}
%respectively.
%\end{defn}
%This is the classical Sugawara construction developed through several papers, including \cite{sugawara, kz}. When acting on insertions, such as
%\[ \lra{\mathfrak{L}_n \phi(x_1) \cdots \phi(x_k)} \]
%only finitely many terms in the series defining $\mathfrak{L}_n$ are non-zero, see Remark~\ref{rem:finitely-many-modes}. The number of non-zero terms depends only on the maximum modulus of $x_1,\ldots,x_k$.
%
%and we require that the same two nested discrete contours separating $0,z_1, \ldots, z_k$ from $\infty$ are used for all pairs $a_{n-j}a_{j}$ and $a_j a_{n-j}$ appearing in \eqref{Ln-definition}. Assuming this, it follows from Lemma~\ref{Dec11.1} that only finitely many terms are non-zero in the infinite sum of \eqref{Ln-definition}. The number of non-zero terms depends on $z_1,\ldots,z_k$ and the choice of contours.
%Informally, we can interpret the $L_n$ as Laurent modes of the stress-energy tensor $T(z) = [J J](z)$, see, e.g., p.47 of \cite{Blumenhagen_Plauschinn}.

\subsection{Virasoro commutation relations}
To be self-contained, and to be
explicit about the well-definedness of the needed compositions of operators,
we include here the well-known calculation of the commutation relations of the
operators $L_n$ obtained by a Sugawara-type construction.
\begin{thm}\label{Ln-GFF}\label{prop:sugawara-construction}\label{thm:gff-virasoro-full-plane}
The two families of change of measure operators $(\bar{L}_n),(L_n), \, n \in \ZZ,$
yield commuting representations of the Virasoro algebra of central charge $c=1$:
\begin{align*}
[L_m, L_n] & \; =  (m-n)L_{m+n} + \frac{1}{12} (m^3-m) \delta_{m+n,0} \rm{Id} , \\
[\bar{L}_m, \bar{L}_n] & \; = (m-n)\bar{L}_{m+n} + \frac{1}{12} (m^3-m) \delta_{m+n,0} \rm{Id} , \\
%\text{and } \qquad
[L_m, \bar{L}_n] & \; = 0.
\end{align*}
We have $[L_n,a_m] = - m a_{n+m}$, $[\bar{L}_n,\bar{a}_m] = - m \bar{a}_{n+m}$,
and  $[L_n,\bar{a}_m] = 0 = [\bar{L}_n, a_m]$.
\end{thm}
%\begin{rem}
%Notice that the theorem gives a precise meaning to the statement that the central charge of the discrete Gaussian free field equals $1$. In particular, we have identified the central charge on the discrete level.
%\end{rem}
\begin{proof}
%[Proof of Theorem~\ref{Ln-GFF}]
By the remark at the end of Section~\ref{sect:current-modes} it is clear that $[L_m, \bar{L}_n] = 0$
for all $m,n \in \mathbb{Z}$. We will check the Virasoro relations for $(L_n)$ by first
treating the actions $(\mathfrak{L}_n)$ on insertions;
the argument for $(\bar{L}_n)$ is identical.
Using the commutation relations of the $\mathfrak{a}_n$ (see the proof of
Theorem~\ref{thm:heisenberg} and Proposition~\ref{prop:current-mode-operators}) and
the simple commutator identity $[A, BC] = [A,B]C + B [A,C]$, we have
\begin{align*}
2[\mathfrak{L}_m, \mathfrak{a}_n]
%= & \, \sum_{j \ge 0} \mathfrak{a}_{m-j}\mathfrak{a}_j \mathfrak{a}_n + \sum_{j \le -1} \mathfrak{a}_j \mathfrak{a}_{m-j} \mathfrak{a}_n - \sum_{j \ge 0} \mathfrak{a}_n \mathfrak{a}_{m-j}\mathfrak{a}_j  - \sum_{j \le -1} \mathfrak{a}_n \mathfrak{a}_j \mathfrak{a}_{m-j} \\
= &  \, \sum_{j \ge 0} \left( [\mathfrak{a}_{m-j}, \mathfrak{a}_n] \mathfrak{a}_j  + \mathfrak{a}_{m-j} [\mathfrak{a}_j, \mathfrak{a}_n] \right) + \sum_{j \le -1} \left( [\mathfrak{a}_j, \mathfrak{a}_n] \mathfrak{a}_{m-j} + \mathfrak{a}_j [ \mathfrak{a}_{m-j} , \mathfrak{a}_n] \right) \\
%= & \, \sum_{j \ge 0} \left( (m-j) \delta_{m-j+n,0} \mathfrak{a}_j +  j \delta_{j+n,0} \mathfrak{a}_{m-j} \right) \\
%  & \qquad + \sum_{j \le -1} \left( j \delta_{j+n,0} \mathfrak{a}_{m-j} + (m-j) \delta_{m-j+n,0} \mathfrak{a}_j \right)\\
= & \, -2 n \mathfrak{a}_{m+n},
\end{align*}
where at each step of the calculation the expression has a well-defined
action on insertions of polynomials $P$ of $\left( \phi(x) \right)_{x \in \mathcal{G}}$
due to a truncation to finitely many terms.
Consequently, again by Theorem~\ref{thm:heisenberg},
Proposition~\ref{prop:current-mode-operators}, and the same commutator identity, we have
\begin{align*}
2[\mathfrak{L}_m, \mathfrak{L}_n] = & \, \sum_{j \ge 0} [\mathfrak{L}_m, \mathfrak{a}_{n-j}\mathfrak{a}_j] + \sum_{j \le -1} [\mathfrak{L}_m, \mathfrak{a}_j\mathfrak{a}_{n-j}]\\
= &  \sum_{j \ge 0}  \left(-(n-j) \mathfrak{a}_{m+n-j} \mathfrak{a}_j - j \mathfrak{a}_{n-j} \mathfrak{a}_{m+j}\right)  \\
& \qquad + \sum_{j \le -1} \left( - (n-j) \mathfrak{a}_j \mathfrak{a}_{m+n-j}-j \mathfrak{a}_{m+j} \mathfrak{a}_{n-j}  \right) \\
= & \sum_{j \ge 0} (m-n+j-m) \mathfrak{a}_{m+n-j} \mathfrak{a}_j  - \sum_{j\ge m} (j-m) \mathfrak{a}_{m+n-j}\mathfrak{a}_j\\
& \qquad + \sum_{j \le -1} (m-n+j-m) \mathfrak{a}_j \mathfrak{a}_{m+n-j} - \sum_{j \le m-1} (j-m) \mathfrak{a}_j \mathfrak{a}_{m+n-j}.
\end{align*}
We can rewrite this last expression as
\begin{align*}
[\mathfrak{L}_m, \mathfrak{L}_n]& =\frac{1}{2}\sum_{j \ge 0} (m-n) \mathfrak{a}_{m+n-j} \mathfrak{a}_j + \frac{1}{2} \sum_{j \le -1} (m-n) \mathfrak{a}_j \mathfrak{a}_{m+n-j} + R_{m,n} \\
& =  (m-n)\mathfrak{L}_{m+n} + R_{m,n},
\end{align*}
where
\[ R_{m,n} := \begin{cases}
    \frac{1}{2}\sum_{j = 0}^{m-1} (m-j)(\mathfrak{a}_j \mathfrak{a}_{m+n-j}- \mathfrak{a}_{m+n-j}\mathfrak{a}_j)
	& \text{if } m>0 \\
    0 & \text{if } m=0 \\
    \frac{1}{2}\sum_{j=m}^{-1} (j-m)(\mathfrak{a}_j \mathfrak{a}_{m+n-j} -\mathfrak{a}_{m+n-j} \mathfrak{a}_j)
	& \text{if } m<0 .
\end{cases} \]
%$R_{0,n} = 0$, and if $m \ge 1$,
%\begin{align*}
%R_{m,n} & =  \frac{1}{2}\sum_{j = 0}^{m-1} (m-j)(\mathfrak{a}_j \mathfrak{a}_{m+n-j}- \mathfrak{a}_{m+n-j}\mathfrak{a}_j) \\
%%& = \frac{1}{2}\sum_{j = 0}^{m-1} (m-j) [\mathfrak{a}_j, \mathfrak{a}_{m+n-j}]  \\
%& =  \frac{1}{2} \sum_{i = 0}^{m-1} (m-j) j \delta_{m+n,0} & \text{(by Proposition~\ref{prop:commutation-relations})} \\
%& =  \frac{1}{12} ( m^3-m) \delta_{m+n,0},
%\end{align*}
We then show that
$ R_{m,n}=\frac{1}{12} ( m^3-m) \delta_{m+n,0}$.
If $m \ge 1$ this is seen by the calculation
\begin{align*}
R_{m,n} & =  \frac{1}{2}\sum_{j = 0}^{m-1} (m-j)(\mathfrak{a}_j \mathfrak{a}_{m+n-j}- \mathfrak{a}_{m+n-j}\mathfrak{a}_j) \\
%& = \frac{1}{2}\sum_{j = 0}^{m-1} (m-j) [\mathfrak{a}_j, \mathfrak{a}_{m+n-j}]  \\
& =  \frac{1}{2} \sum_{j = 0}^{m-1} (m-j) j \delta_{m+n,0} & \text{(by Proposition~\ref{prop:commutation-relations})} \\
& =  \frac{1}{12} ( m^3-m) \delta_{m+n,0} ,
\end{align*}
and if $m \le -1$, the calculation is similar.

As in the proof of Theorem~\ref{thm:heisenberg}, the commutation
relations for the actions on insertions give the commutation relations for the change of
measure operators, because $L_n a_m - a_m L_n$ is the operator obtained by lifting the action
of $\mathfrak{a}_{-m} \mathfrak{L}_{-n} - \mathfrak{L}_{-n} \mathfrak{a}_{-m}$ in the
sense of Lemma~\ref{lem: lifting to change of measure operators}, and
$L_n L_m - L_m L_n$ is the operator obtained by lifting the action
of $\mathfrak{L}_{-m} \mathfrak{L}_{-n} - \mathfrak{L}_{-n} \mathfrak{L}_{-m}$ etc.
\end{proof}

%//////////////////////////////

%\section{Half plane}
%In this section we will sketch a construction of $\tilde{L}_n$:s for the discrete GFF in the upper half-plane. Let $\tilde{\phi}
%(z), z \in \ZZ^2_\diamond \cap \mathbb{H}$ be the discrete Gaussian free field with Dirichlet boundary condition on $\ZZ^2_\diamond \cap \mathbb{R}$. (It is convenient to set $\tilde \phi \equiv 0$ on the lower half-plane.) The two-point correlation function for $\tilde\phi$ is
%\[
%\lra{\tilde{\phi}(z) \tilde{\phi}(w)}_{\mathbb{H}} = \tilde{G}(z,w)\]
%where
%\begin{equation}\label{GH}
%\tilde{G}(z,w) : = G(z,w) - G(z,\bar{w}), \quad z,w \in \ZZ^2_\diamond,
%\end{equation}
%and $G(z,w)$ is the free Green's function for the discrete Laplacian on $\ZZ^2_\diamond$. From this it follows that for $z,w \in \ZZ^2_\diamond$,
%\begin{align}\label{GHrels}
%\tilde{G}(z,w) & = \tilde{G}(\bar{z},\bar{w})=-\tilde{G}(\bar{z},w)  =- \tilde{G}(z,\bar{w}).
%\end{align}
%When $z \in \ZZ^2_m$ is in $\mathbb{H}$ we have
%\[
%\lra{\dd\tilde{\phi}(z) \tilde{\phi}(w)}_\HH= \dd_1 \tilde{G}(z,w), \quad \lra{\ddbar \tilde{\phi}(z) \tilde{\phi(w)}}_\HH = \ddbar_1 \tilde{G}(z,w) = - \dd_1 \tilde{G}(\bar{z},w).
%\]
%while if $x \in \ZZ^2_m$ is real then since $\tilde\phi$ is $0$ on the lower half-plane we can write
%\begin{align}\label{Greal}
%\lra{\dd \tilde{\phi}(x) \tilde{\phi}(w)}_{\HH} & = -\frac{i}{2}\tilde{G}\left(x+\frac{i}{2},w\right) \\
%& =\frac{1}{2}\dd_1 \tilde{G}(x, w)\end{align}
%and
%\begin{align}\label{Greal2}
%\lra{\ddbar \tilde{\phi}(x) \tilde{\phi}(w)}_{\HH} & = \frac{i}{2}\tilde{G}\left(x+\frac{i}{2},w\right) \\
%& =-\frac{1}{2}\dd_1 \tilde{G}(x, w).
%\end{align}



\section{Further results}
\subsection{Coulomb gas}\label{sect:coulomb-gas}
By adding a multiple of the current mode to each $L_n$ one obtains representations of
the Virasoro algebra with other central charges. For $b \in \bR $ we define change of
measure operators
\begin{align}\label{def:cg}
L_n^b:=L_n+b(n+1) a_n, \qquad \bar{L}_n^b:=\bar{L}_n+b(n+1) \bar{a}_n,
\end{align}
which will both satisfy the Virasoro commutation relations with central charge $c=1-12b^2$.
For $b \neq 0$, these operators are not parity preserving (nor are they parity reversing),
as opposed to the $L_n$ corresponding to $c=1$, and they can therefore be considered less
natural.
%It is not clear if
%one can construct the $L_n^b$ by considering current modes of a modified
%discrete field.
\begin{prop}\label{prop:coulomb_gas}
For $c \leq 1$ and $n\in\mathbb{Z}$, if we define the change
of measure operators $L_{n}^{b},\bar{L}_{n}^{b} \colon \mathcal{M}\to\mathcal{M}$
by \eqref{def:cg} with $b=\pm\sqrt{\frac{1-c}{12}}$, we have
that $\left(L_{n}^{b}\right)_{n\in\mathbb{Z}}$
and $\left(\bar{L}_{n}^{b}\right)_{n\in\mathbb{Z}}$ yield two commuting
representations of the Virasoro algebra with central charge $c$, that is
$[L_n^b , \bar{L}_m^b] = 0$,
\[ [L_m^b, L_n^b]
    = (m-n)L^b_{m+n} + \frac{c}{12}(m^3-m) \delta_{m+n,0} \, \id_\mathcal{M}, \]
and similarly for $(\bar{L}_n^b)$.
\end{prop}
\begin{proof}
Starting from the observation  $[L_m^b, a_n] = -n a_{m+n} + b m \delta_{m+n,0} \id_\mathcal{M}$,
the proof is a straightforward modification of that of Theorem~\ref{prop:sugawara-construction}.
\end{proof}


\subsection{Discrete GFF in a half plane}
A similar construction based on current modes can be carried out for the discrete GFF
in the upper half-plane $\bH = \bZ \times \bN$ with Dirichlet boundary
condition on $\partial \bH = \ZZ \times \set{0}$. We sketch the construction.

In the half-plane $\bH$, the Green's function $G^\bH(z,w)$ is well defined
%as the expected number of visits to $w$ of a simple random walk from $z$ in $\bH$,
even without an infrared regularization: for a fixed $w \in \bH$
it is determined by $- [\Delta G^\bH](\cdot,w) = \delta_w$ and
$G^\bH(z,w)=0$ for $z \in \partial \bH$.
%and it is given by %$G^{\bH}(z , w) = \potker(z-\bar{w}) - \potker(z-w)$.
%\[ G^{\bH}(x + \ii y , x' + \ii y') =
%\frac{1}{(2\pi)^2} \iint_{[-\pi,\pi]^2}
%    \frac{e^{\ii (x-x') q + \ii (y-y') p} - e^{\ii (x-x') q + \ii (y+y') p}}
%        {1 - \frac{1}{2} \cos(q) - \frac{1}{2} \cos(p)} \ud q \ud p .\]
The discrete GFF in the upper half-plane is a collection $(\phi^{\bH}(z))_{z \in \bH}$
of centered real Gaussians with covariance given by the Green's function
\[ \EX\left[ \phi^{\bH}(z) \phi^{\bH}(w) \right] = G^{\bH}(z , w) .\]

We note that since $G^{\bH}(z , w) = \potker(z-\bar{w}) - \potker(z-w)$, the field
$\phi^\bH$ can be constructed by a reflection: if $\phi=(\phi(z))_{z \in \bZ^2}$ is
the pinned discrete GFF on the infinite square lattice $\bZ^2$, then the field
\[ \phi^\bH (z) = \frac{1}{\sqrt{2}} \Big( \phi(z) - \phi(\bar{z}) \Big) \]
is a discrete GFF in the upper half-plane. We define the current by
\[ J^\bH (z) = [\dd \phi^{\bH}](z) \]
with $\phi^\bH$ extended to the dual as zero. With the reflection trick, the current
makes sense for all $z \in \bZ^2_m$, its correlation functions are recovered by Wick's
formula from the two-point functions
\[ \lra{J^\bH (z) \phi^\bH(x)} = \Cker(z-x) - \Cker(z-\bar{x}) \]
for $x \in \bZ^2$ (discrete holomorphic for $z\notin\set{x,\bar{x}}$),
and we can define the action on insertions by integration
\[ \lra{\mathfrak{a}^\bH_n f(\phi^\bH \big|_{\mathcal{G}}) }_\mu :=
\frac{1}{\sqrt{\pi}} \oint_{[\gamma]} \lra{ J^\bH(z) f(\phi^\bH \big|_{\mathcal{G}})g(\phi^\bH \big|_{\mathcal{G}})}
    z^{[n]} \, \ud z \]
where $[\gamma]$ is any sufficiently large simple closed discrete contour
and $g=g^\mu_{\mathcal{G}}$ is a Radon-Nikodym derivative.
If one prefers not to resort to the reflection trick, the above expression can
be written entirely in terms of $(\phi^{\bH}(z))_{z \in \bH}$ by noting that
the contour may be assumed to be symmetric with respect to a reflection in the real axis,
and then in the integrand we may use
$\bar{z}^{[n]}=\overline{z^{[n]}}$ and $ J^\bH(\bar{z}) = - \bar{J}^\bH(z)$ to rewrite
everything in terms of only the part of the contour that is in the upper half plane.

From the two-point function of the current
\[ \lra{J^\bH (z) J^\bH (w)} = [\dd \Cker](w-z) - \delta_{z,\bar{w}} \]
one recovers the Heisenberg commutation relations
\[ \mathfrak{a}^\bH_m \mathfrak{a}^\bH_n - \mathfrak{a}^\bH_n \mathfrak{a}^\bH_m = m \delta_{n+m,0} \, \id \]
as before. The construction of the change of measure operators $(a^\bH_n)_{n\in\bZ}$ and the
Virasoro operators $(L^\bH_n)_{n\in\bZ}$ is entirely parallel to the case treated earlier, and
one obtains the following result:

\begin{thm}\label{thm:half-plane}
Let $\mathcal{M}^\bH$ be the vector space of complex Gibbs measures that are
changes of measure of law of the discrete GFF in $\bH$, and whose
Radon-Nikodym derivatives with respect to it are in $L^p$ for all $p<\infty$.

There exists a family of change of measure operators $(a^\bH_n)_{n\in\bZ}$
on $\mathcal{M}^\bH$ that are parity reversing and that satisfy the Heisenberg
algebra commutation relations %$[a^\bH_m,a^\bH_n]=m \delta_{n+m,0} \id$
and a family of change of measure operators $(L^\bH_n)_{n\in\bZ}$ on $\mathcal{M}^\bH$
that are parity preserving and that satisfy the Virasoro algebra
commutation relations with central charge $c=1$ %, $[a^\bH_m,a^\bH_n]=m \delta_{n+m,0} \id$
and such that $[L_n,a_m] = - m a_{n+m}$.
\end{thm}



%\begin{align*}
%%\lra{\phi(x_1) \cdots \phi(x_k)}_{\mathbb{H}} \mapsto
%\frac{1}{i}\int_{[\sigma]} \lra{ [\dd \phi ](z) \phi(x_1) \cdots \phi(x_k)}_{\mathbb{H}} z^{[n]} \, dz + \frac{1}{i}\int_{[\sigma]}\lra{[ \ddbar \phi] \phi(x_1) \cdots \phi(x_k)}_{\mathbb{H}} \bar{z}^{[n]}\, d\bar{z}.
%\end{align*}
%Here $[\sigma]$ is a discrete ``crosscut'' defined by the disjoint union of the edges of $\dd_0 \cE_{\Sigma}$ and $\dd_1 \cE_{\Sigma_m}$ that intersect $\HH$ for some discrete domain $\Sigma$ such that $\partial \Sigma \cap \HH$ is a crosscut of $\HH$ separating $0, x_1, \ldots x_k$ from $\infty$ in $\HH$.
%
%The $\tilde{\mathfrak{a}}_n$ lift to change of measure operators and yield a representation of the Heisenberg algebra: using the reflection symmetry of the correlation function of the half-plane discrete GFF one can show that the two integrals defining the $\tilde{\mathfrak{a}}_n$ combine and can be written as one standard discrete contour integral involving symmetrized versions of the whole-plane correlation kernels. The Virasoro modes can then be defined by the Sugawara construction.

%\subsection{\red{Ising model}}

\bibliographystyle{9}

\begin{thebibliography}{References}

%\bibitem[AsMc11]{assis-mccoy}M. Assis, B. M. McCoy, The energy density
%of an Ising half-plane lattice, \emph{J. Phys. A:Math. Theor.} 44:095003,
%2011.
%[Bax89]

\bibitem[BaBe03]{bauer-bernard-i}M.~Bauer, D.~Bernard, Conformal field theories of stochastic {L}oewner evolutions, \emph{Comm. Math. Phys.} 239: 493--521, 2003.

\bibitem[BaBe04]{bauer-bernard-ii}M.~Bauer, D.~Bernard, Conformal transformations and the SLE partition function martingale, \emph{Ann. Henri Poincar\'e} 5, 289, 2004. % [arXiv:math-ph/0305061]

\bibitem[Bax89]{baxter}R. Baxter, \emph{Exactly solved models in
statistical mechanics}. Academic Press Inc. {[}Harcourt Brace Jovanovich
Publishers{]}, London, 1989.

\bibitem[BeHo13]{benoist-hongler} S. Benoist, C. Hongler, In preparation, 2013.

\bibitem[BPZ84a]{belavin-polyakov-zamolodchikov-i}A. A. Belavin,
A. M. Polyakov, A. B. Zamolodchikov. Infinite conformal symmetry in
two-dimensional quantum field theory\emph{. Nucl. Phys. B}, 241(2):333--380,
1984.

\bibitem[BPZ84b]{belavin-polyakov-zamolodchikov-ii}A. A. Belavin,
A. M. Polyakov, A. B. Zamolodchikov. Infinite conformal symmetry of
critical fluctuations in two dimensions.\emph{ J. Stat. Phys.}, 34(5-6):763--774,
1984.

%\bibitem[BoDT09]{boutillier-de-tiliere-i}C. Boutillier, B. de Tili?re,
%The critical Z-invariant Ising model via dimers: locality property.
%\emph{Comm. Math. Phys}, to appear. arXiv:0902.1882v1, 2009.
%
%\bibitem[BoDT08]{boutillier-de-tiliere-ii}C. Boutillier, B. de Tili?re,
%The critical Z-invariant Ising model via dimers: the periodic case.
%\emph{PTRF} 147:379-413, 2010. arXiv:0812.3848v1.

%\bibitem[BuGu93]{burkhard-guim}T. Burkhardt, I. Guim, \emph{Conformal
%theory of the two-dimensional Ising model with homogeneous boundary
%conditions and with disordered boundary fields}, Phys. Rev. B (1),
%47:14306-14311, 1993.

%\bibitem[CaNe07a]{camia-newman-i}F. Camia and C.M. Newman, Critical
%percolation exploration path and SLE(6): a proof of convergence\emph{.
%Probab. Th. Rel. Fields} 139:473-519, 2007.

%\bibitem[CaNe07b]{camia-newman-ii}F. Camia and C. M. Newman,  Two-Dimensional
%Critical Percolation: The Full Scaling Limit\emph{. Comm. Math. Phys.}
%268(1):1--38, 2007.\emph{ }

\bibitem[Car84]{cardy-i}J. Cardy, Conformal invariance and Surface
Critical Behavior\emph{. Nucl. Phys. B} 240:514--532, 1984.

%\bibitem[Car92]{cardy-iii}J. Cardy. Critical percolation in finite
%geometries\emph{. J. Phys. A} 25(4):L201\textendash{}L206, 1992.

%\bibitem[ChIz11]{chelkak-izyurov}D. Chelkak and K. Izyurov, Holomorphic
%Spinor Observables in the Critical Ising Model, arXiv:1105.5709v2.

\bibitem[CHI12]{chi-ising}D.~Chelkak, C.~Hongler, K. Izyurov, Conformal Invariance of Spin Correlations in the Planar Ising Model, arXiv:1202.2838, 2012.

\bibitem[ChSm11]{chelkak-smirnov-i}D. Chelkak and S. Smirnov, Discrete
complex analysis on isoradial graphs\emph{. Advances in Mathematics},
228:1590--1630, 2011.

\bibitem[ChSm09]{chelkak-smirnov-ii}D. Chelkak and S. Smirnov, Universality
in the 2D Ising model and conformal invariance of fermionic observables\emph{.
Inventiones} \emph{Math.}, to appear. arXiv:0910.2045.

%
%\bibitem[ChSm11]{chelkak-smirnov-iii}D. Chelkak and S. Smirnov,\emph{
%}preprint.

%\bibitem[DeMo10]{dembo-montanary}A. Dembo, A. Montanari, Ising model
%on locally tree-like graphs, \emph{Ann. Appl. Probab.}, 20(2):565--592,
%2010.

\bibitem[DMS97]{di-francesco-mathieu-senechal}P. Di Francesco, P.
Mathieu, D. S\'en\'echal, \emph{Conformal Field Theory}. Graduate texts
in contemporary physics. Springer-Verlag New York, 1997.

\bibitem[DRC06]{cardy-doyon-riva} B.~Doyon, V.~Riva, J.~Cardy, Identification of the stress-energy tensor through conformal restriction in {SLE} and related processes. \emph{Comm. Math. Phys.},268:687--716, 2006.

\bibitem[Dub13]{dubedat-def} J. Dub\'edat, personal communication.

\bibitem[Duf56]{duffin}R.~J.~Duffin, Basic properties of discrete analytic functions. \emph{Duke Math. J.}, 23: 335--363, 1956.

\bibitem[FrWe03]{werner-friedrich-i}R. Friedrich, W. Werner, Conformal restriction, highest-weight representations and SLE, \emph{Comm. Math. Phys.}, 243: 105--122, 2003.

\bibitem[Fun05]{funaki}T.~Funaki, Stochastic interface models, \emph{Lecture Notes in Math.}, 1869: {103--274}, 2005.

%\bibitem[Dub11a]{dubedat-i}J. Dub?dat, Dimers and analytic torsion
%I, arXiv:1110.2808v1.

%\bibitem[Dub11b]{dubedat-ii}J. Dub?dat\emph{, }Exact bosonization
%of the Ising model, arXiv:1112.4399v1.
%
%\bibitem[Dub12]{dubedat-iii}J. Dub?dat, in preparation.
%
%\bibitem[DHN11]{duminil-copin-hongler-nolin}H. Duminil-Copin, C.
%Hongler, P. Nolin, Connection probabilities and RSW-type bounds for
%the FK Ising model\emph{. Comm. Pure. Appl. Math. }64(9):1165--1198,
%2011\emph{.}

%\bibitem[Fis98]{fisher}M. E. Fisher, Renormalization group theory:
%its basis and formulation in statistical mechanics. \emph{Rev. Modern
%Phys.}, 70(2):653--681, 1998.

\bibitem[Gaw99]{gawedzki}
K. Gaw{\c e}dzki,
Lectures on conformal field theory.
In \emph{Quantum fields and strings: A course for
mathematicians}, Vols. 1, 2 (Princeton, NJ, 1996/1997), pages 727--805.
Amer. Math. Soc., Providence, RI, 1999.

%\bibitem[GRS08]{garban-rohde-schramm}C. Garban, S. Rohde, O. Schramm,
%Continuity of the SLE trace in simply connected domains. \emph{Israel
%Journal of Math., }to appear, 2008.

%\bibitem[GHS70]{griffith-hurst-sherman}R. Griffiths, C.A. Hurst,
%S. Sherman, Concavitiy of Magnetization of an Ising Ferromagnet in
%a Positive External Field, \emph{J. Math. Phys. }11(3):790--795, 1970.

%\bibitem[Gri06]{grimmett}G. Grimmett, \emph{The Random-Cluster Model}.
%Volume 333 of Grundlehren der Mathematischen Wissenschaften, Springer-Verlag,
%Berlin, 2006.

%\bibitem[Hec67]{hecht}R. Hecht, Correlation Functions for the Two-Dimensional
%Ising Model, \emph{Phys. Rev.} 158:557--561, 1967.

\bibitem[Hon10]{hongler-i}C. Hongler, \emph{Conformal invariance
of Ising model correlations}. Ph.D. thesis, University of Geneva,
http://www.math.columbia.edu/\textasciitilde{}hongler/thesis.pdf,
2010.


%\bibitem[HoKy13]{hongler-kytola} C. Hongler, K. Kyt\"ol\"a, In preparation, 2013.

%\bibitem[HoSm10]{hongler-smirnov-ii}C. Hongler, S. Smirnov.\emph{
%}The energy density in the critical planar Ising model. arXiv:1008.2645v3.
%
%\bibitem[Isi25]{ising}E. Ising, Beitrag zur Theorie des Ferromagnetismus\emph{.
%Zeitschrift f?r Physik}, 31:253--258, 1925.

\bibitem[ItTh87]{itoyama-thacker}H.~Itoyama, H.~B.~Thacker, Lattice Virasoro algebra and corner transfer matrices in the Baxter eight-vertex model. \emph{Phys. Rev. Lett.} 58: 1395--1398, 1987.

\bibitem[Jan97]{janson}S.~Janson, \emph{Gaussian Hilbert spaces.} Cambridge Tracts in Mathematics 129, Cambridge University Press 1997.

\bibitem[KaMa11]{kang-makarov}N.-G.~Kang, N. G.~Makarov. Gaussian free field and conformal field theory. arXiv:1101.1024, 2011.
%
%\bibitem[KaCe71]{kadanoff-ceva}L. Kadanoff, H. Ceva, Determination
%of an operator algebra for the two-dimensional Ising model\emph{.
%Phys. Rev. B} (3), 3:3918--3939, 1971.

\bibitem[Ken00]{kenyon-ci}R.~Kenyon, Conformal invariance of domino tiling. \emph{Ann. Probab.} 28(2): 759--795, 2000.

\bibitem[KoSa94]{koo-saleur} W.M. Koo, H. Saleur, Representations of the Virasoro algebra from lattice models, \emph{Nucl. Phys. B}, 426:459--504, 1994.

\bibitem[KnZa84]{kz}V.~G.~Knizhnik, A.~B.~Zamoldchikov, Current algebra and Wess-Zumino model in two dimensions. \emph{Nucl. Phys. B} 247: 83--103, 1984.

%\bibitem[KrWa41]{kramers-wannier}H. A. Kramers and G. H. Wannnier,
%Statistics of the two-dimensional ferromagnet. I.\emph{ Phys. Rev.
%(2)}, 60:252--262, 1941.

\bibitem[Kyt07]{kytola-local_mgales_virasoro} K.~Kyt\"ol\"a, Virasoro module structure of local martingales of SLE variants, \emph{Rev. Math. Phys.}
19, 455, 2007. % [arXiv:math-ph/0604047]

%\bibitem[LaLi04]{lawler-limi}G.F. Lawler, V. Limic, The Beurling
%Estimate for a Class of Random Walks, \emph{Elec. J. Prob. 9(27)},
%2004.

\bibitem[LaLi10]{lawler-limic}G. F.~Lawler, V.~Limic, \emph{Random walk: a modern introduction.} Cambridge.

\bibitem[LSW04]{lsw-lerw}G. F.~Lawler, O.~Schramm,W.~Werner, Conformal invariance of planar loop-erased random walks and uniform spanning trees. \emph{Ann. Probab.}, 32(1b): 939--995, 2004.
%
%\bibitem[Len20]{lenz}W. Lenz,\emph{ }Beitrag zum Verst?ndnis der
%magnetischen Eigenschaften in festen K?rpern. \emph{Phys. Zeitschr.},
%21:613--615, 1920.

\bibitem[JM94]{jimbo-miwa}M. Jimbo, T. Miwa, \emph{Algebraic Analysis of Solvable Lattice Models}, CBMS Regional Conferenfece Series in Mathematics, 85, 1994.

\bibitem[Mer01]{mercat}C.~Mercat, Discrete Riemann surfaces and the Ising model. \emph{Comm. Math. Phys.}, 218(1):177--216, 2001.

\bibitem[Mic89]{mickelsson}J. Mickelsson, \emph{Current Algebras and Groups}, Plenum Monographs in Nonlinear Physics, Springer, 1989.

\bibitem[Sch00]{schramm-sle} O. Schramm, Scaling limits of loop-erased random walks and uniform spanning trees. \emph{Israel J. Math.} 118: 221--288, 2000.

\bibitem[ScSh09]{schramm-sheffield-gff-i} O.~Schramm, S.~Sheffield, Contour lines of the two-dimensional discrete {G}aussian free field. \emph{Acta Math.} 202: 21--137, 2009.

\bibitem[She05]{sheffield-gff-ptrf}S.~Sheffield, Gaussian free fields for mathematicians. \emph{Probab. Theory Related Fields}. 139: 521--541, 2007.

\bibitem[Smi01]{smirnov-percolation}S.~Smirnov, Critical percolation in the plane: conformal invariance, Cardy's formula, scaling limits. \emph{C. R. Acad. Sci. Paris Ser. I Math.}, 333(3): 239--244, 2001.

\bibitem[Smi10]{smirnov-ising}S.~Smirnov, Conformal invariance in random cluster models. I. Holomorphic fermions in the Ising model. \emph{Ann. of Math.}172(2): 1435--1467, 2010.
\bibitem[Sug68]{sugawara}H.~Sugawara, A field theory of currents. \emph{Phys. Rev. Lett.} 170:1659, 1968.
\bibitem[Som68]{sommerfield}C. Sommerfield, Currents as dynamical variables, \emph{Phys. Rev. Lett.} 176:2019, 1968.
%\bibitem[McWu67]{mccoy-wu-ii}B. M. McCoy and T. T. Wu, Theory of
%Toeplitz Determinants and the Spin Correlations of the Two-Dimensional
%Ising Model. IV, \emph{Phys. Rev. }162(2):436--475, 1967.

%\bibitem[McWu73]{mccoy-wu-i}B. M. McCoy and T. T. Wu, \emph{The two-dimensional
%Ising model}. Harvard University Press, Cambridge, Massachusetts,
%1973.

%\bibitem[Ons44]{onsager}L. Onsager, Crystal statistics. I. A two-dimensional
%model with an order-disorder transition\emph{. Phys. Rev. }65(2):117--149,
%1944.
%
%\bibitem[Pal07]{palmer}J. Palmer, \emph{Planar Ising correlations}.
%Birkh?user, 2007.
%
%\bibitem[PaTr83]{palmer-tracy}J. Palmer, C. A. Tracy,\emph{ }Two-Dimensional
%Ising Correlations: The SMJ Analysis.\emph{ Adv. in Appl. Math. }4:46--102,
%1983.
%
%\bibitem[Pin11]{pinson}H. Pinson, Rotational Invariance of the 2D
%spin-spin correlation function.
%
%\bibitem[SMJ77]{sato-miwa-jimbo-i}R. Peierls, On Ising's model of
%ferromagnetism. \emph{Proc. Camb. Philos. Soc., }32:477--481, 1936.M.
%Sato, T. Miwa, M. Jimbo,\emph{ }Studies on holonomic quantum fields,
%I\emph{-IV. Proc. Japan Acad. Ser. A Math. Sci., }53(1):6--10, 53(1):147-152,
%53(1):153-158, 53(1):183-185, 1977.
%
%\bibitem[SMJ79a]{sato-miwa-jimbo-ii}M. Sato, T. Miwa, M. Jimbo, Holonomic
%quantum fields III.\emph{ Publ. RIMS, Kyoto Univ. }15:577--629, 1979.
%
%\bibitem[SMJ79b]{sato-miwa-jimbo-iii}M. Sato, T. Miwa, M. Jimbo,
%Holonomic quantum fields IV.\emph{ Publ. RIMS, Kyoto Univ.} 15:871--972,
%1979.
%
%\bibitem[SMJ80]{sato-miwa-jimbo-iv}M. Sato, T. Miwa, M. Jimbo,\emph{
%}Holonomic quantum fields V. \emph{Publ. RIMS, Kyoto Univ.} 16:531--584,
%1980.
%
%\bibitem[Smi06]{smirnov-i}S. Smirnov,\emph{ }Towards conformal invariance
%of 2D lattice models\emph{.} \emph{Sanz-Sol?, Marta (ed.) et al.,
%Proceedings of the international congress of mathematicians (ICM),
%Madrid, Spain, August 22--30, 2006. Volume II: Invited lectures, 1421-1451.
%Z?rich: European Mathematical Society (EMS),} 2006.
%
%\bibitem[WMTB76]{wu-mccoy-tracy-barouch}T. T. Wu, B. M. McCoy, C.
%A. Tracy, E. Barouch, Spin-spin correlation functions for the two-dimensional
%Ising model: Exact theory in the scaling region.\emph{ Phys. Rev.
%B }13:316--374, 1976.

\end{thebibliography}

\end{document}
